%=============================================================
%  QUANTUM MECHANICS & QUANTUM CHEMISTRY — COURSE SYLLABUS
%  Restructured hierarchy: Series > Module > Lecture
%=============================================================
\documentclass[11pt,a4paper]{article}

%--- Packages ------------------------------------------------
\usepackage[T1]{fontenc}
\usepackage[utf8]{inputenc}
\usepackage{fix-cm}
\usepackage[activate={false},final]{microtype}
\usepackage[margin=2.5cm]{geometry}
\usepackage{amsmath,amssymb,amsthm}
\usepackage{physics}
\usepackage{braket}
\usepackage{booktabs}
\usepackage{array}
\usepackage{longtable}
\usepackage{tabularx}
\usepackage{multirow}
\usepackage{enumitem}
\usepackage{sectsty}
\usepackage{fancyhdr}
\usepackage[dvipsnames,svgnames,table]{xcolor}
\usepackage{hyperref}
\usepackage{parskip}
\usepackage{tcolorbox}
\tcbuselibrary{skins,breakable}

%--- Colours -------------------------------------------------
\definecolor{darkblue}{RGB}{10,40,90}
\definecolor{midblue}{RGB}{30,80,160}
\definecolor{lightblue}{RGB}{220,233,255}
\definecolor{accentgold}{RGB}{180,140,20}
\definecolor{lightgold}{RGB}{255,248,220}
\definecolor{darkgray}{RGB}{60,60,60}
\definecolor{lightgray}{RGB}{245,245,245}
\definecolor{chemgreen}{RGB}{20,100,60}
\definecolor{lightgreen}{RGB}{220,245,230}

%--- Section fonts via sectsty (no titlesec = no conflict) ---
\sectionfont{\Large\bfseries\color{darkblue}\sectionrule{0ex}{0pt}{-1ex}{1.5pt}}
\subsectionfont{\large\bfseries\color{midblue}}
\subsubsectionfont{\normalsize\bfseries\color{darkgray}}

%--- Hyperref ------------------------------------------------
\hypersetup{
  colorlinks  = true,
  linkcolor   = midblue,
  urlcolor    = midblue,
  citecolor   = midblue,
  pdftitle    = {Quantum Mechanics and Quantum Chemistry -- Course Syllabi},
  pdfauthor   = {},
}

%--- Header / Footer -----------------------------------------
\pagestyle{fancy}
\fancyhf{}
\fancyhead[L]{\small\color{darkgray}\textit{Quantum Mechanics \& Quantum Chemistry}}
\fancyhead[R]{\small\color{darkgray}\thepage}
\fancyfoot[C]{\small\color{darkgray}\textit{For academic use only}}
\renewcommand{\headrulewidth}{0.4pt}

%--- tcolorbox styles ----------------------------------------
\tcbset{
  modulebox/.style={
    enhanced, breakable,
    colback=lightblue, colframe=midblue,
    fonttitle=\bfseries\color{white}, coltitle=white,
    attach boxed title to top left={yshift=-2mm, xshift=4mm},
    boxed title style={colback=midblue, colframe=darkblue,
      sharp corners, boxrule=0.5pt},
    sharp corners=south,
    boxrule=0.8pt,
    left=6pt, right=6pt, top=8pt, bottom=6pt,
  },
  chembox/.style={
    enhanced, breakable,
    colback=lightgreen, colframe=chemgreen,
    fonttitle=\bfseries\color{white}, coltitle=white,
    attach boxed title to top left={yshift=-2mm, xshift=4mm},
    boxed title style={colback=chemgreen, colframe=chemgreen!80!black,
      sharp corners, boxrule=0.5pt},
    sharp corners=south,
    boxrule=0.8pt,
    left=6pt, right=6pt, top=8pt, bottom=6pt,
  },
  highlightbox/.style={
    enhanced,
    colback=lightgold, colframe=accentgold,
    boxrule=0.6pt,
    left=5pt, right=5pt, top=4pt, bottom=4pt,
    fontupper=\small\itshape,
  },
  seriesbox/.style={
    enhanced, breakable,
    colback=lightgold, colframe=accentgold,
    fonttitle=\bfseries\color{darkblue}, coltitle=darkblue,
    attach boxed title to top left={yshift=-2mm, xshift=4mm},
    boxed title style={colback=lightgold, colframe=accentgold,
      sharp corners, boxrule=0.5pt},
    sharp corners=south,
    boxrule=1pt,
    left=6pt, right=6pt, top=8pt, bottom=6pt,
  },
}

\newcommand{\coursebadge}[1]{%
  \begin{tcolorbox}[highlightbox]#1\end{tcolorbox}%
}

%=============================================================
\begin{document}
%=============================================================

%--- Title Page ----------------------------------------------
\begin{titlepage}
\pagecolor{darkblue}
\color{white}
\vspace*{3cm}
\begin{center}
  {\fontsize{28}{34}\selectfont\bfseries Quantum Mechanics}\\[0.4cm]
  {\fontsize{18}{24}\selectfont\bfseries \& Quantum Chemistry}\\[0.6cm]
  {\color{accentgold}\rule{10cm}{2pt}}\\[0.6cm]
  {\large\itshape Comprehensive Course Syllabi}\\[0.3cm]
  {\normalsize Physics \& Theoretical Chemistry Programme}\\[2.5cm]
  \begin{tabular}{ll}
    \textbf{Series~I:}   & Introductory Quantum Mechanics\\
                         & \quad Module~I.1 --- State Space, Operators, and the Bra-Ket Language (6 lectures)\\
                         & \quad Module~I.2 --- Bra-Ket Notation and the Language of Quantum Mechanics (12 lectures)\\
                         & \quad Module~I.3 --- Quantum Dynamics, Perturbation Theory, and Physical Systems (11 lectures)\\[6pt]
    \textbf{Series~II:}  & Advanced Quantum Mechanics\\
                         & \quad 7 Modules (II.1--II.7), 12 lectures each\\[6pt]
    \textbf{Series~III:} & Molecular Electronic Structure \& Quantum Chemistry\\
                         & \quad 3 Modules (III.1--III.3), approx.\ 36--40 lectures total\\
  \end{tabular}\\[2.5cm]
  {\color{accentgold}\rule{10cm}{1pt}}\\[0.4cm]
  {\small\itshape Each module in Series~I \& II comprises 12 lectures unless otherwise noted.\\
  Hierarchy: \textbf{Series} $\supset$ \textbf{Module} $\supset$ \textbf{Lecture}}
\end{center}
\vfill
\end{titlepage}
\nopagecolor

%--- Table of Contents ---------------------------------------
\tableofcontents
\newpage

%=============================================================
\section*{Preamble and Programme Overview}
\addcontentsline{toc}{section}{Preamble and Programme Overview}
%=============================================================

\coursebadge{%
  This document describes a sequential, research-oriented programme in quantum
  mechanics and quantum chemistry, organised into a three-level hierarchy:
  \textbf{Series} (broad thematic strand) $\supset$ \textbf{Module} (self-contained
  teaching unit, typically 12 lectures) $\supset$ \textbf{Lecture} (single 90-minute
  session).\\[4pt]
  \textbf{Series~I} introduces the algebraic and dynamical foundations of quantum
  mechanics across three modules:
  Module~I.1 introduces the Hilbert-space language of operators and states;
  Module~I.2 develops the full bra-ket formalism via a dual Dirac--Sakurai approach
  with five-tier pedagogy;
  Module~I.3 covers quantum dynamics, all pictures of quantum mechanics, perturbation
  theory, open systems, and example physical systems.\\[2pt]
  \textbf{Series~II} deepens the formalism across seven thematic modules.
  \textbf{Series~III} pivots to molecular electronic structure and computational
  quantum chemistry.
}

\medskip
\noindent\textbf{Notation conventions used throughout:}
\begin{itemize}[leftmargin=1.5em]
  \item $\hat{A}, \hat{H}, \hat{U}$ --- operators (hat notation)
  \item $\ket{\psi}, \bra{\phi}$ --- Dirac ket/bra states
  \item $\hat{\rho}$ --- density operator/matrix
  \item $[\hat{A},\hat{B}]=\hat{A}\hat{B}-\hat{B}\hat{A}$ --- commutator
  \item $\{\hat{A},\hat{B}\}=\hat{A}\hat{B}+\hat{B}\hat{A}$ --- anti-commutator
\end{itemize}

\medskip
\noindent\textbf{Hierarchy key:}
\begin{itemize}[leftmargin=1.5em]
  \item \textbf{Series} --- top-level strand labelled by Roman numeral (I, II, III)
  \item \textbf{Module} --- self-contained teaching block labelled \textit{Series.Number}
    (e.g.\ I.1, II.4, III.2)
  \item \textbf{Lecture} --- individual 90-min session labelled L\textit{NN} within its
    module (e.g.\ L01, L07, L12)
\end{itemize}

\newpage

%#############################################################
%  SERIES I SEPARATOR PAGE
%#############################################################
\newpage
\thispagestyle{empty}
\pagecolor{darkblue}
\color{white}
\vspace*{2.2cm}
\begin{center}
  {\color{accentgold}\rule{11cm}{1.5pt}}\\[0.5cm]
  {\fontsize{11}{14}\selectfont\bfseries\color{accentgold}\MakeUppercase{Series I}}\\[0.3cm]
  {\fontsize{26}{32}\selectfont\bfseries Introductory Quantum Mechanics}\\[0.4cm]
  {\color{accentgold}\rule{11cm}{1.5pt}}\\[1.0cm]
  {\large\itshape Three modules --- 29 lectures $\times$ 90 minutes}\\[0.5cm]
  {\normalsize\itshape Audience: Advanced undergraduates through early graduate students}\\[0.3cm]
  {\normalsize\itshape Prerequisites: Linear algebra, classical mechanics, basic complex analysis}\\[1.6cm]

  \begin{tabular}{@{}clr@{}}
    \toprule
    \textbf{Module} & \textbf{Title} & \textbf{Lectures}\\
    \midrule
    \textbf{I.1} & State Space, Operators, and the Bra-Ket Language & 6\\[3pt]
                 & \multicolumn{2}{@{}l}{\small\color{accentgold}%
                   Hilbert space; bounded/unbounded operators; adjoint; unitary maps;}\\
                 & \multicolumn{2}{@{}l}{\small\color{accentgold}%
                   Hermitian observables; spectral theorem; operator algebra;}\\
                 & \multicolumn{2}{@{}l}{\small\color{accentgold}%
                   exponential map; BCH lemma; Wigner's theorem; bra-ket formalism.}\\[6pt]
    \textbf{I.2} & Bra-Ket Notation and the Language of Quantum Mechanics & 12\\[3pt]
                 & \multicolumn{2}{@{}l}{\small\color{accentgold}%
                   States as rays; dual space; inner product; linear operators;}\\
                 & \multicolumn{2}{@{}l}{\small\color{accentgold}%
                   eigenkets; projectors; completeness; continuous bases;}\\
                 & \multicolumn{2}{@{}l}{\small\color{accentgold}%
                   measurement theory; uncertainty; tensor products; spin; density operators;}\\
                 & \multicolumn{2}{@{}l}{\small\color{accentgold}%
                   rigged Hilbert spaces. Dual Dirac--Sakurai approach; five-tier pedagogy.}\\[6pt]
    \textbf{I.3} & Quantum Dynamics, Perturbation Theory, and Physical Systems & 11\\[3pt]
                 & \multicolumn{2}{@{}l}{\small\color{accentgold}%
                   Density matrix; discrete/continuous bases; uncertainty; symmetries;}\\
                 & \multicolumn{2}{@{}l}{\small\color{accentgold}%
                   Noether's theorem; time evolution (I \& II); Schr\"odinger, Heisenberg,}\\
                 & \multicolumn{2}{@{}l}{\small\color{accentgold}%
                   and Dirac pictures; perturbation theory; open systems; example systems.}\\[4pt]
    \midrule
    & \textbf{Total} & \textbf{29}\\
    \bottomrule
  \end{tabular}

  \vfill
  {\color{accentgold}\rule{11cm}{0.8pt}}\\[0.4cm]
  {\small\itshape\color{white!80!darkblue}%
    Each lecture: 90 minutes $\cdot$ Five-tier pedagogy (HS through PhD) $\cdot$
    Dual Dirac--Sakurai approach}\\[0.3cm]
  {\color{accentgold}\rule{11cm}{0.8pt}}
\end{center}
\newpage
\nopagecolor
\color{black}

%#############################################################
\section{Series I --- Introductory Quantum Mechanics}
\label{sec:seriesI}
%#############################################################

\begin{tcolorbox}[seriesbox, title={Series~I Overview}]
\textbf{Audience:} Advanced undergraduates through early graduate students.\\
\textbf{Prerequisites:} Linear algebra, classical mechanics, basic complex analysis.\\
\textbf{Structure:} Three modules totalling 29 lectures $\times$ 90~minutes.\\[4pt]
\begin{tabular}{llr}
\toprule
\textbf{Module} & \textbf{Title} & \textbf{Lectures}\\
\midrule
I.1 & State Space, Operators, and the Bra-Ket Language & 6\\
I.2 & Bra-Ket Notation and the Language of Quantum Mechanics & 12\\
I.3 & Quantum Dynamics, Perturbation Theory, and Physical Systems & 11\\
\midrule
 & \textbf{Total} & \textbf{29}\\
\bottomrule
\end{tabular}
\end{tcolorbox}

%=============================================================
\subsection{Module I.1 --- State Space, Operators, and the Bra-Ket Language}
\label{sec:mod-I1}
%=============================================================

\coursebadge{%
  \textbf{Intended audience:} Advanced undergraduates with a solid background in
  classical mechanics and linear algebra.\\[2pt]
  \textbf{Scope:} Foundational module covering the algebraic language of quantum
  mechanics --- from the structure of Hilbert space and operator theory through
  to the introduction of quantum states, Dirac notation, and the density matrix
  formalism up to and including the notion of pure and mixed states (\S2.3).\\[2pt]
  \textbf{Duration:} 6 lectures $\times$ 90 minutes (Lectures~L01--L06,
  sections 2.1--2.3). Module~I.2 develops the full bra-ket formalism in
  depth; Module~I.3 continues from \S2.4 onward.
}

\subsubsection*{Module Overview}
\addcontentsline{toc}{subsubsection}{Module I.1 Overview}

%------------------------------------------------------------
\subsubsection*{Module Overview}
\addcontentsline{toc}{subsubsection}{Module I.1 Overview}

\begin{center}
\begin{tabularx}{\textwidth}{clX}
\toprule
\textbf{Lec.} & \textbf{Title} & \textbf{Key Concepts}\\
\midrule
L01 & The Hilbert Space of Quantum States &
    Hilbert space $\mathcal{H}$; vectors, norms, completeness; bounded vs.\ unbounded operators; adjoint\\
L02 & Unitary Operators and Symmetry Transformations &
    $\hat{U}\hat{U}^\dagger=\hat{1}$; inner-product preservation; rotations, translations, time evolution\\
L03 & Observables, Hermitian Operators, and the Spectral Theorem &
    Observables $\leftrightarrow$ Hermitian operators; real eigenvalues; orthogonal eigenstates; spectral decomposition\\
L04 & Operator Algebra and the Exponential Map &
    Sums, products, functions of operators; $e^{i\theta\hat{A}}$; Baker--Campbell--Hausdorff lemma; commutators\\
L05 & Types of Symmetry Operators and Wigner's Theorem &
    Unitary vs.\ anti-unitary symmetries; Wigner's theorem; projectors $\hat{P}^2=\hat{P}$; time-reversal\\
L06 & Bra-Ket Formalism (\S2.1--2.2) &
    Dirac notation; ket/bra duality; inner product; outer product; matrix elements; expectation values\\
\bottomrule
\end{tabularx}
\end{center}

%------------------------------------------------------------
\subsubsection*{Five-Tier Layered Pedagogy}
\addcontentsline{toc}{subsubsection}{Module I.1 Five-Tier Pedagogy}

Each lecture is stratified across five audience levels, enabling delivery to a mixed cohort:

\begin{center}
\begin{tabular}{llp{9cm}}
\toprule
\textbf{Tier} & \textbf{Cohort} & \textbf{Characteristic content}\\
\midrule
Tier~1 & HS       & Geometric intuition, analogy with familiar vector spaces, conceptual motivation\\
Tier~2 & Beg.\ UG & Definitions, worked matrix examples, basic proofs; algebraic facility\\
Tier~3 & Adv.\ UG & Rigorous derivations, Cauchy--Schwarz, spectral theorem, proof habits\\
Tier~4 & MSc      & Operator domains, subtle distinctions (Hermitian vs.\ self-adjoint), advanced applications\\
Tier~5 & PhD      & Functional analysis, deficiency indices, unbounded operators, mathematical physics depth\\
\bottomrule
\end{tabular}
\end{center}

%------------------------------------------------------------
\subsubsection*{Lecture-by-Lecture Syllabi}
\addcontentsline{toc}{subsubsection}{Module I.1 Lecture Syllabi}

% ------ L01 -----------------------------------------------
\begin{tcolorbox}[modulebox, title={Module~I.1 --- Lecture L01: The Hilbert Space of Quantum States}]
\textbf{90-min pacing:} (0--20~min) Motivation --- why ordinary Euclidean space is insufficient;
infinite-dimensional state spaces and physical necessity;
(20--45~min) Definition of a Hilbert space $\mathcal{H}$: inner product, norm, completeness (Cauchy sequences);
(45--70~min) Bounded vs.\ unbounded operators; operator norm; adjoint $\hat{A}^\dagger$ defined via
$\langle\phi|\hat{A}\psi\rangle = \langle\hat{A}^\dagger\phi|\psi\rangle$;
(70--90~min) Anti-Hermitian operators; $[\hat{A}^\dagger]^\dagger=\hat{A}$; worked examples.

\medskip\textbf{Core learning outcomes:}
\begin{itemize}[leftmargin=1.4em,topsep=2pt,itemsep=1pt]
  \item State the definition of a Hilbert space and verify that $L^2(\mathbb{R})$ is one.
  \item Distinguish bounded from unbounded operators and give physical examples of each.
  \item Define the adjoint $\hat{A}^\dagger$ and compute it for $2\times 2$ matrices.
  \item Classify Hermitian ($\hat{A}=\hat{A}^\dagger$) and anti-Hermitian ($\hat{A}=-\hat{A}^\dagger$) operators.
\end{itemize}

\medskip\textbf{Dual-track content:}
\begin{itemize}[leftmargin=1.4em,topsep=2pt,itemsep=1pt]
  \item \textit{Geometric track:} Hilbert space as the ``right'' infinite-dimensional generalisation
    of $\mathbb{R}^n$; inner product $\langle\cdot|\cdot\rangle$ and angles; completeness as
    ``no missing limits.''
  \item \textit{Algebraic/axiomatic track:} Formal axioms (conjugate symmetry, linearity, positive
    definiteness); operator theory; norm from inner product $\|\psi\|=\sqrt{\langle\psi|\psi\rangle}$.
\end{itemize}

\medskip\textbf{Tier differentiation:}
HS: inner product as ``overlap''; norm as ``length''; $\mathbb{R}^2$ and $\mathbb{C}^2$ as tangible examples.
Beg.\ UG: verify axioms in $\mathbb{C}^n$; compute $\hat{A}^\dagger$ for explicit matrices.
Adv.\ UG: prove adjoint is unique; show $(\hat{A}\hat{B})^\dagger = \hat{B}^\dagger\hat{A}^\dagger$.
MSc: $L^2$ space rigorously; operator norm; boundedness and continuity equivalence.
PhD: spectrum of unbounded operators; domain issues; closed operators; Hellinger--Toeplitz theorem.

\medskip\textbf{Assessment:}
Set~A --- verify Hilbert space axioms for $\mathbb{C}^n$; compute adjoints of given $3\times 3$ matrices.
Set~B --- prove $(\hat{A}^\dagger)^\dagger = \hat{A}$; show bounded $\Rightarrow$ continuous.
Projects: (simple) norm visualiser for $L^2$ functions; (complex) counterexample: unbounded operator with dense domain.
\end{tcolorbox}

\medskip

% ------ L02 -----------------------------------------------
\begin{tcolorbox}[modulebox, title={Module~I.1 --- Lecture L02: Unitary Operators and Symmetry Transformations}]
\textbf{90-min pacing:} (0--15~min) Recap: inner products and adjoint;
(15--40~min) Definition of unitary operator: $\hat{U}^\dagger\hat{U} = \hat{U}\hat{U}^\dagger = \hat{1}$;
preservation of inner products and norms $\langle\hat{U}\phi|\hat{U}\psi\rangle = \langle\phi|\psi\rangle$;
(40--65~min) Physical incarnations: spatial rotations $\hat{U}(\theta)$, spatial translations $\hat{T}(a)$,
time-evolution operator $\hat{U}(t) = e^{-i\hat{H}t/\hbar}$ (preview);
(65--90~min) Unitary equivalence of operators; spectral invariance under unitary transformations.

\medskip\textbf{Core learning outcomes:}
\begin{itemize}[leftmargin=1.4em,topsep=2pt,itemsep=1pt]
  \item Verify unitarity $\hat{U}^\dagger\hat{U}=\hat{1}$ from definition and examples.
  \item Prove that unitary maps preserve norms and transition probabilities.
  \item Identify unitary operators arising from physical symmetries (rotation, translation, phase).
  \item Understand $\hat{U} = e^{i\hat{G}}$ for Hermitian generator $\hat{G}$ (Stone's theorem preview).
\end{itemize}

\medskip\textbf{Dual-track content:}
\begin{itemize}[leftmargin=1.4em,topsep=2pt,itemsep=1pt]
  \item \textit{Physical track:} Symmetry $\to$ conservation law (Noether preview); rotations as $SU(2)$
    action; phase freedom $e^{i\theta}\hat{1}$ as trivial unitary.
  \item \textit{Mathematical track:} Unitary group $\mathcal{U}(\mathcal{H})$; isometric isomorphisms;
    one-parameter unitary groups $\hat{U}(t)$; Stone's theorem (statement).
\end{itemize}

\medskip\textbf{Tier differentiation:}
HS: rotations as length-preserving maps; ``rotating'' a quantum state without losing information.
Beg.\ UG: verify $\hat{U}^\dagger\hat{U}=\hat{1}$ for rotation matrices; compute $\hat{U}\ket{\psi}$.
Adv.\ UG: show eigenvalues of unitary operators lie on the unit circle; prove spectral invariance.
MSc: one-parameter groups; Stone's theorem; generator $\hat{G}$ and Lie algebra connection.
PhD: strongly continuous unitary groups on Banach spaces; Hille--Yosida theorem; domain of generator.

\medskip\textbf{Assessment:}
Set~A --- verify unitarity for $2\times 2$ rotation and phase matrices; compute $\hat{U}\ket{\psi}$.
Set~B --- prove unitary $\Rightarrow$ isometric; show spectrum of $\hat{U}$ lies on $|z|=1$.
Projects: (simple) animate unitary evolution of a two-level state; (complex) matrix exponential $e^{i\theta\hat{G}}$ for various generators.
\end{tcolorbox}

\medskip

% ------ L03 -----------------------------------------------
\begin{tcolorbox}[modulebox, title={Module~I.1 --- Lecture L03: Observables, Hermitian Operators, and the Spectral Theorem}]
\textbf{90-min pacing:} (0--20~min) The measurement postulate: physical observables correspond to Hermitian operators;
motivation from real-valued experimental outcomes;
(20--45~min) Key properties of Hermitian operators: real eigenvalues, orthogonality of eigenvectors
($\langle a|a'\rangle = \delta_{aa'}$), completeness of eigenbasis;
(45--70~min) The spectral theorem (discrete case): $\hat{A} = \sum_a a\ket{a}\bra{a}$;
functions of Hermitian operators $f(\hat{A})$;
(70--90~min) Degenerate eigenvalues; continuous spectrum preview; physical examples
($\hat{x}$, $\hat{p}$, $\hat{H}$ for a particle in a box).

\medskip\textbf{Core learning outcomes:}
\begin{itemize}[leftmargin=1.4em,topsep=2pt,itemsep=1pt]
  \item State and apply the measurement postulate linking observables to Hermitian operators.
  \item Prove that Hermitian operators have real eigenvalues and orthogonal eigenvectors.
  \item Write the spectral decomposition $\hat{A}=\sum_a a\ket{a}\bra{a}$ and apply it.
  \item Compute $f(\hat{A})$ from the spectral decomposition for simple functions.
\end{itemize}

\medskip\textbf{Dual-track content:}
\begin{itemize}[leftmargin=1.4em,topsep=2pt,itemsep=1pt]
  \item \textit{Physical track (Sakurai):} Preparation and measurement; Stern--Gerlach as paradigm;
    eigenvalues as the only possible measurement outcomes.
  \item \textit{Abstract track (Dirac):} Spectral theorem as the operational backbone; resolution of
    identity $\hat{1}=\sum_a\ket{a}\bra{a}$; postulate as \emph{definition} of observable.
\end{itemize}

\medskip\textbf{Tier differentiation:}
HS: ``measuring an observable always gives one of its eigenvalues'' --- Stern--Gerlach conceptual story.
Beg.\ UG: find eigenvalues/vectors of $2\times 2$ Hermitian matrices; write spectral form.
Adv.\ UG: prove real eigenvalues and orthogonality; handle degenerate case with Gram--Schmidt.
MSc: spectral theorem in full generality; functions of operators; self-adjoint vs.\ Hermitian distinction preview.
PhD: spectral measure and projection-valued measure (PVM); Borel functional calculus; continuous spectrum treatment.

\medskip\textbf{Assessment:}
Set~A --- diagonalise Hermitian matrices; compute $e^{i\hat{A}}$ via spectral decomposition.
Set~B --- prove spectral theorem for $n\times n$ Hermitian matrices; degenerate subspace orthogonalisation.
Projects: (simple) eigenvalue/eigenvector visualiser for $2\times 2$ observables; (complex) spectral resolution
of the harmonic oscillator Hamiltonian.
\end{tcolorbox}

\medskip

% ------ L04 -----------------------------------------------
\begin{tcolorbox}[modulebox, title={Module~I.1 --- Lecture L04: Operator Algebra and the Exponential Map}]
\textbf{90-min pacing:} (0--15~min) Operator arithmetic: sums $(\hat{A}+\hat{B})\ket{\psi}$,
products $\hat{A}\hat{B}\ket{\psi}$, scalar multiples; non-commutativity;
(15--40~min) Functions of operators via power series and spectral decomposition;
the matrix exponential $e^{\hat{A}} = \sum_{n=0}^\infty \hat{A}^n/n!$ and its convergence;
(40--65~min) The exponential map $e^{i\theta\hat{A}}$ for Hermitian $\hat{A}$: unitarity proof;
one-parameter families; physical applications (phase, rotation, evolution operators);
(65--90~min) Baker--Campbell--Hausdorff (BCH) lemma: $e^{\hat{A}}e^{\hat{B}} = e^{\hat{A}+\hat{B}+\frac{1}{2}[\hat{A},\hat{B}]+\cdots}$;
statement, significance, and low-order applications.

\medskip\textbf{Core learning outcomes:}
\begin{itemize}[leftmargin=1.4em,topsep=2pt,itemsep=1pt]
  \item Perform operator arithmetic and recognise non-commutativity.
  \item Compute matrix exponentials via diagonalisation and power series.
  \item Prove $\hat{U}=e^{i\theta\hat{A}}$ is unitary when $\hat{A}$ is Hermitian.
  \item State the BCH lemma and apply it to first order in the commutator.
\end{itemize}

\medskip\textbf{Dual-track content:}
\begin{itemize}[leftmargin=1.4em,topsep=2pt,itemsep=1pt]
  \item \textit{Algebraic track:} Operator ring structure; Lie brackets and Lie algebras;
    BCH as the structure theorem linking Lie algebra to Lie group.
  \item \textit{Computational track:} Step-by-step matrix exponential evaluation;
    $e^{i\theta\hat{\sigma}_z}$ and $e^{i\phi\hat{\sigma}_x}$ as worked examples; physical interpretation.
\end{itemize}

\medskip\textbf{Tier differentiation:}
HS: ``exponential of a matrix'' --- analogy with $e^x$ as infinite sum; rotation by small angles.
Beg.\ UG: compute $e^{i\theta\hat{\sigma}_z}$; verify unitarity numerically.
Adv.\ UG: prove unitarity of $e^{i\theta\hat{A}}$ rigorously; derive BCH to first order.
MSc: full BCH series; Zassenhaus formula; group--algebra correspondence; Magnus expansion preview.
PhD: convergence of BCH in Banach algebras; analytic extensions; applications in quantum control.

\medskip\textbf{Assessment:}
Set~A --- compute $e^{i\theta\hat{\sigma}_z}$, $e^{i\phi\hat{\sigma}_y}$; verify unitarity; simple BCH applications.
Set~B --- prove BCH to second order; derive Hadamard identity $e^{\hat{A}}\hat{B}e^{-\hat{A}}$.
Projects: (simple) matrix exponential calculator with visualisation; (complex) first-order BCH vs.\ exact
product comparison for non-commuting $2\times 2$ Hamiltonians.
\end{tcolorbox}

\medskip

% ------ L05 -----------------------------------------------
\begin{tcolorbox}[modulebox, title={Module~I.1 --- Lecture L05: Types of Symmetry Operators and Wigner's Theorem}]
\textbf{90-min pacing:} (0--20~min) Taxonomy of operators: linear vs.\ anti-linear;
unitary vs.\ anti-unitary; physical motivation from symmetries that reverse arrow of time;
(20--50~min) Wigner's theorem: every quantum symmetry is either a unitary or anti-unitary operator
on $\mathcal{H}$; statement, intuition, and significance;
(50--70~min) Projectors $\hat{P} = \ket{\phi}\bra{\phi}$: idempotency $\hat{P}^2=\hat{P}$,
Hermiticity, eigenvalues $\{0,1\}$; orthogonal projectors and decomposition of identity;
(70--90~min) Time-reversal as paradigmatic anti-unitary symmetry; Kramers' theorem preview.

\medskip\textbf{Core learning outcomes:}
\begin{itemize}[leftmargin=1.4em,topsep=2pt,itemsep=1pt]
  \item Distinguish linear/unitary from anti-linear/anti-unitary operators with examples.
  \item State Wigner's theorem and explain its physical necessity.
  \item Construct rank-1 projectors $\hat{P}=\ket{\phi}\bra{\phi}$ and verify $\hat{P}^2=\hat{P}$.
  \item Explain how time-reversal is modelled by an anti-unitary operator.
\end{itemize}

\medskip\textbf{Dual-track content:}
\begin{itemize}[leftmargin=1.4em,topsep=2pt,itemsep=1pt]
  \item \textit{Physical track:} Symmetry as a map on quantum states preserving transition probabilities;
    Wigner's theorem as the mathematical conclusion; time-reversal and Kramers degeneracy in spin systems.
  \item \textit{Mathematical track:} Anti-linear maps $\mathcal{H}\to\mathcal{H}$;
    anti-unitary condition $\langle\Theta\phi|\Theta\psi\rangle = \langle\psi|\phi\rangle$;
    projector algebra and projection theorem.
\end{itemize}

\medskip\textbf{Tier differentiation:}
HS: symmetry as ``doing something and the physics looks the same''; time-reversal intuitively.
Beg.\ UG: compute $\hat{P}^2$ for rank-1 projectors; verify $\hat{P}$ Hermitian.
Adv.\ UG: prove eigenvalues of $\hat{P}$ are $0$ and $1$; orthogonal projector resolution of identity.
MSc: Wigner's theorem proof sketch; complex conjugation as anti-unitary; Kramers' theorem statement.
PhD: full proof of Wigner's theorem via Uhlhorn's theorem; anti-unitary representations; spin-statistics connection preview.

\medskip\textbf{Assessment:}
Set~A --- verify $\hat{P}^2=\hat{P}$ for explicit projectors; identify unitary vs.\ anti-unitary from action on basis.
Set~B --- prove eigenvalues of a projector are $\{0,1\}$; show time-reversal squaring gives $\pm\hat{1}$ for integer/half-integer spin.
Projects: (simple) symmetry classification quiz for common quantum operators;
(complex) Kramers degeneracy verification for spin-$\tfrac{1}{2}$ Hamiltonian under time-reversal.
\end{tcolorbox}

\medskip

% ------ L06 (former L02) -----------------------------------
\begin{tcolorbox}[modulebox, title={Module~I.1 --- Lecture L06: Bra-Ket Formalism (\S2.1--2.2)}]
\textbf{90-min pacing:} (0--20~min) From column vectors to abstract kets:
$\ket{\psi}\in\mathcal{H}$; superposition and linearity;
(20--45~min) The dual space $\mathcal{H}^*$; bra $\bra{\phi}$ as linear functional;
inner product $\braket{\phi|\psi}$; conjugate symmetry;
(45--70~min) Outer product $\ket{\psi}\bra{\phi}$; matrix element $\bra{\phi}\hat{A}\ket{\psi}$;
expectation value $\langle\hat{A}\rangle = \bra{\psi}\hat{A}\ket{\psi}$; variance $(\Delta A)^2$;
(70--90~min) Notation dictionary: abstract ket $\leftrightarrow$ column vector $\leftrightarrow$ wavefunction;
connection to previous lectures; bridge to Module~I.2.

\medskip\textbf{Core learning outcomes:}
\begin{itemize}[leftmargin=1.4em,topsep=2pt,itemsep=1pt]
  \item Write kets and bras in Dirac notation and perform inner/outer products.
  \item Evaluate matrix elements $\bra{\phi}\hat{A}\ket{\psi}$ for operators from Lectures L01--L05.
  \item Compute expectation values and variances of observables.
  \item Translate freely between abstract Dirac notation, column-vector representation, and wavefunction language.
\end{itemize}

\medskip\textbf{Dual-track content:}
\begin{itemize}[leftmargin=1.4em,topsep=2pt,itemsep=1pt]
  \item \textit{Sakurai (operational) track:} Notation as a concise bookkeeping device;
    amplitude $\braket{a|\psi}$ as measurement probability amplitude; $|\braket{a|\psi}|^2$ as Born rule.
  \item \textit{Dirac (abstract) track:} Bra as element of algebraic dual space $\mathcal{H}^*$;
    Riesz representation theorem identifies $\bra{\psi}$ with $\ket{\psi}$ via inner product.
\end{itemize}

\medskip\textbf{Tier differentiation:}
HS: ``bra'' and ``ket'' as ``row'' and ``column'' in a richer language; inner product as overlap.
Beg.\ UG: compute $\braket{\phi|\psi}$, outer products, and $\langle\hat{A}\rangle$ for $2\times 2$ examples.
Adv.\ UG: prove $(\ket{\psi}\bra{\phi})^\dagger = \ket{\phi}\bra{\psi}$; work with sums of projectors.
MSc: Riesz representation theorem (statement and significance); dual space in infinite dimensions.
PhD: bounded vs.\ unbounded functionals; Dirac notation in the context of rigged Hilbert spaces.

\medskip\textbf{Assessment:}
Set~A --- inner products, outer products, expectation values; verify Hermiticity of $\hat{x}$ in position representation.
Set~B --- prove variance formula $(\Delta A)^2 = \langle\hat{A}^2\rangle - \langle\hat{A}\rangle^2$;
show $(\hat{A}\ket{\psi}\bra{\phi})^\dagger = \ket{\phi}\bra{\psi}\hat{A}^\dagger$.
Projects: (simple) Dirac notation calculator for spin-$\tfrac{1}{2}$; (complex) derive Born-rule consequences
from inner-product axioms and completeness.
\end{tcolorbox}

\textit{Continued in Module~I.3, Lecture~L01: pure and mixed states,
density matrix properties, and partial trace (\S2.3--2.5).}

\newpage
%=============================================================
\subsection{Module I.2 --- Bra-Ket Notation and the Language of Quantum Mechanics}
\label{sec:mod-I2}
%=============================================================

\coursebadge{%
  \textbf{Module designation:} Series~I, Module~2 --- develops the full bra-ket
  formalism in depth, bridging the algebraic foundations of Module~I.1 with the
  dynamical and structural content of Module~I.3.\\[2pt]
  \textbf{Dual-track approach:} Dirac (abstract axiomatic) + Sakurai (operational/physical),
  with a \textbf{five-tier layered pedagogy} spanning HS through PhD.\\[2pt]
  \textbf{Duration:} 12 lectures $\times$ 90 minutes. Spiral curriculum: seven competency
  threads revisited progressively across lectures.\\[2pt]
  \textbf{Assessment:} Per-lecture bundles (concept check, Problem Sets A \& B, mini-projects,
  research questions) plus course-level grading. See Section~\ref{sec:braket-assessment}.
}

\subsubsection*{Module Overview}
\addcontentsline{toc}{subsubsection}{Module I.2 Overview}

\begin{center}
\begin{tabularx}{\textwidth}{clX}
\toprule
\textbf{Lec.} & \textbf{Title} & \textbf{Key Content}\\
\midrule
L01 & Why Dirac Notation? States as Rays \& Superposition &
      States, kets, rays, global phase, Born rule preview; Sakurai operational vs.\ Dirac abstract\\
L02 & Bras, Dual Space \& Inner Product Structure &
      Dual vectors, conjugate linearity, inner product axioms, Riesz theorem (intuition)\\
L03 & Linear Operators: Adjoint, Hermitian \& Unitary &
      Operator algebra, matrix elements $\bra{\phi}\hat{A}\ket{\psi}$, observables \& symmetries\\
L04 & Eigenkets, Projectors \& Spectral Decomposition &
      Measurement outcomes, projectors $\hat{P}_a=\ket{a}\bra{a}$, post-measurement state\\
L05 & Completeness, Change of Basis \& Matrix Mechanics &
      Identity insertion, basis transforms, matrix representation $A_{mn}=\bra{m}\hat{A}\ket{n}$\\
L06 & Continuous Bases: $\ket{x}$, $\ket{p}$ \& Delta Functions &
      Position/momentum bases, wavefunctions, Fourier connection, distributions\\
L07 & Measurement Theory: Postulates, Compatibility \& Uncertainty &
      Commutators, Robertson--Schr\"odinger inequality, compatible observables, POVM preview\\
L08 & Time Evolution \& Pictures of Quantum Mechanics &
      Unitary evolution, Schr\"odinger/Heisenberg pictures, constants of motion, Ehrenfest\\
L09 & Tensor Products: Composite Systems \& Entanglement &
      Product vs.\ entangled states, Schmidt decomposition, partial measurement\\
L10 & Spin \& Finite-Dimensional Hilbert Spaces ($SU(2)$) &
      Spin-$\tfrac{1}{2}$, Pauli matrices, rotation operators, Bloch sphere, $SU(2)$ vs.\ $SO(3)$\\
L11 & Density Operators, Mixed States \& Partial Trace &
      $\hat{\rho}=\sum_i p_i\ket{\psi_i}\bra{\psi_i}$, $\operatorname{Tr}(\hat{\rho}\hat{A})$,
      partial trace, Bloch sphere, purity\\
L12 & Rigged Hilbert Spaces \& Mathematical Foundations &
      Gel'fand triple $\Phi\subset\mathcal{H}\subset\Phi'$, distributions, self-adjointness
      vs.\ Hermiticity, safe usage rules\\
\bottomrule
\end{tabularx}
\end{center}

%------------------------------------------------------------
\subsubsection*{Spiral Curriculum --- Competency Threads}
\addcontentsline{toc}{subsubsection}{Module I.2 Spiral Curriculum}

The module is built around seven interlocking threads, each revisited at increasing depth
across the 12 lectures:

\begin{center}
\begin{tabular}{lp{10cm}}
\toprule
\textbf{Thread} & \textbf{Content arc}\\
\midrule
\textbf{States \& Duality} & Kets/bras, inner products, amplitudes (L01--L02)\\
\textbf{Representation} & Basis expansion, discrete/continuous, Fourier (L03, L05--L06)\\
\textbf{Operators} & Adjoint, Hermitian/unitary, spectral theorem (L03--L04)\\
\textbf{Measurement} & Projectors, collapse, POVM preview (L04, L07)\\
\textbf{Dynamics} & Unitary evolution, Schr\"odinger/Heisenberg pictures (L08)\\
\textbf{Composition} & Tensor products, entanglement, partial trace (L09, L11)\\
\textbf{Foundations} & Distributions, operator domains, rigged Hilbert spaces (L06, L12)\\
\bottomrule
\end{tabular}
\end{center}

%------------------------------------------------------------
\subsubsection*{Five-Tier Layered Pedagogy}
\addcontentsline{toc}{subsubsection}{Module I.2 Five-Tier Pedagogy}

Each lecture is explicitly stratified across five audience levels, enabling delivery
to a mixed cohort:

\begin{center}
\begin{tabular}{llp{9cm}}
\toprule
\textbf{Tier} & \textbf{Cohort} & \textbf{Characteristic content}\\
\midrule
Tier~1 & HS       & Complex numbers, amplitudes, vector addition analogy,
                    intuitive motivation\\
Tier~2 & Beg.\ UG & Normalization, orthogonality, components $c_n=\braket{n|\psi}$,
                    computational facility\\
Tier~3 & Adv.\ UG & Cauchy--Schwarz, spectral theorem, proof habits, rigorous derivations\\
Tier~4 & MSc      & Projective Hilbert space, density operators, measurement subtleties,
                    representation theory\\
Tier~5 & PhD      & Rigged Hilbert space, unbounded operators, functional analysis,
                    deficiency indices\\
\bottomrule
\end{tabular}
\end{center}

%------------------------------------------------------------
\subsubsection*{Lecture-by-Lecture Syllabi}
\addcontentsline{toc}{subsubsection}{Module I.2 Lecture Syllabi}

\begin{tcolorbox}[modulebox, title={Module~I.2 --- Lecture L01: Why Dirac Notation? States as Rays \& Superposition}]
\textbf{90-min pacing:} (0--15~min) Motivation --- representations vs.\ abstract states;
(15--45~min) Kets as vectors, rays, normalisation, global phase;
(45--70~min) Bras (informal), amplitudes, Born rule preview;
(70--90~min) Notation dictionary: ket $\leftrightarrow$ column vector $\leftrightarrow$ wavefunction.

\medskip\textbf{Core learning outcomes:}
\begin{itemize}[leftmargin=1.4em,topsep=2pt,itemsep=1pt]
  \item Understand state as equivalence class $\ket{\psi}\sim e^{i\theta}\ket{\psi}$; global phase unobservable.
  \item Interpret amplitudes and Born-rule probabilities $P(a)=|\braket{a|\psi}|^2$.
  \item Build the notation dictionary in a discrete basis.
  \item Connect Sakurai (operational) and Dirac (abstract) perspectives.
\end{itemize}

\medskip\textbf{Dual-track content:}
\begin{itemize}[leftmargin=1.4em,topsep=2pt,itemsep=1pt]
  \item \textit{Sakurai track:} Preparation procedures $\to$ stable frequencies; interference motivates
    complex amplitudes; $P=|A_1+A_2|^2$.
  \item \textit{Dirac track:} $\ket{\psi}$ as ray in $\mathcal{H}$; completeness
    $\sum_n\ket{n}\bra{n}=\hat{1}$; $P(a_i)=|\braket{a_i|\psi}|^2$.
\end{itemize}

\medskip\textbf{Tier differentiation:}
HS: $|z|^2$ as probability, interference analogy.
Beg.\ UG: normalisation $\braket{\psi|\psi}=1$, components.
Adv.\ UG: Cauchy--Schwarz, triangle inequality.
MSc: rays and projective Hilbert space $\mathbb{P}(\mathcal{H})$, Wigner preview.
PhD: separability, precise Hilbert space statements.

\medskip\textbf{Assessment:}
Set~A --- normalisation + amplitude problems.
Set~B --- prove Cauchy--Schwarz; phase-invariance exercises.
Projects: (simple) basis-expansion calculator; (complex) compare physical meaning across representations.
\end{tcolorbox}

\medskip
\begin{tcolorbox}[modulebox, title={Module~I.2 --- Lecture L02: Bras, Dual Space \& Inner Product Structure}]
\textbf{90-min pacing:} (0--20~min) Define bras as linear functionals on $\mathcal{H}$;
(20--50~min) Conjugate linearity; inner product axioms;
(50--70~min) Probability amplitude $\braket{\phi|\psi}$ and Born rule;
(70--90~min) Orthonormal basis; Gram--Schmidt sketch.

\medskip\textbf{Core learning outcomes:}
\begin{itemize}[leftmargin=1.4em,topsep=2pt,itemsep=1pt]
  \item Define bra $\bra{\psi}$ as element of dual space $\mathcal{H}^*$.
  \item Master conjugate linearity and inner product rules.
  \item Derive Born-rule probability from inner product structure.
  \item Riesz representation theorem (intuition).
\end{itemize}

\medskip\textbf{Tier differentiation:}
HS: conjugation rules, dot-product analogy.
Beg.\ UG: orthonormal bases, Gram--Schmidt.
Adv.\ UG: Riesz representation theorem intuition, proofs.
MSc: anti-linear maps, adjoint rigorously, continuity.
PhD: bounded vs.\ unbounded functionals, Banach space context.

\medskip\textbf{Assessment:}
Set~A --- compute overlaps, orthogonality checks.
Set~B --- inequalities, norm properties, inner-product geometry.
Projects: (simple) visualise state geometry; (complex) derive Born-rule consequences from axioms.
\end{tcolorbox}

\medskip
\begin{tcolorbox}[modulebox, title={Module~I.2 --- Lecture L03: Linear Operators: Adjoint, Hermitian \& Unitary}]
\textbf{90-min pacing:} (0--20~min) Operators as maps $\mathcal{H}\to\mathcal{H}$; matrix elements;
(20--50~min) Adjoint $\hat{A}^\dagger$; Hermitian observables; unitary symmetries;
(50--70~min) Expectation values; operator algebra;
(70--90~min) Commutators preview; spectral theorem motivation.

\medskip\textbf{Core learning outcomes:}
\begin{itemize}[leftmargin=1.4em,topsep=2pt,itemsep=1pt]
  \item Compute matrix elements $\bra{\phi}\hat{A}\ket{\psi}$.
  \item Define adjoint: $\bra{\phi}\hat{A}\ket{\psi}=\bra{\hat{A}^\dagger\phi}\ket{\psi}$.
  \item Characterise Hermitian ($\hat{A}=\hat{A}^\dagger$) and unitary ($\hat{U}^\dagger\hat{U}=\hat{1}$) operators.
  \item Prove expectation value of Hermitian operator is real.
\end{itemize}

\medskip\textbf{Tier differentiation:}
HS: matrices as machines on vectors, $2\times2$ examples.
Beg.\ UG: eigenvalues, $\langle\hat{A}\rangle=\bra{\psi}\hat{A}\ket{\psi}$.
Adv.\ UG: proof Hermitian $\Rightarrow$ real expectation; unitary preserves inner product.
MSc: spectral theorem preview, commutators, normal operators.
PhD: domains of unbounded operators, self-adjoint vs.\ merely Hermitian.

\medskip\textbf{Assessment:}
Set~A --- compute adjoint, check Hermiticity, expectation values.
Set~B --- prove spectral properties; operator algebra.
Projects: (simple) classify $2\times 2$ matrices; (complex) commutator algebra and uncertainty preview.
\end{tcolorbox}

\medskip
\begin{tcolorbox}[modulebox, title={Module~I.2 --- Lecture L04: Eigenkets, Projectors \& Spectral Decomposition}]
\textbf{90-min pacing:} (0--20~min) Eigenvalue equation $\hat{A}\ket{a}=a\ket{a}$;
(20--50~min) Projectors $\hat{P}_a=\ket{a}\bra{a}$; properties;
(50--75~min) Spectral decomposition $\hat{A}=\sum_a a\ket{a}\bra{a}$;
(75--90~min) Measurement postulate; post-measurement state.

\medskip\textbf{Core learning outcomes:}
\begin{itemize}[leftmargin=1.4em,topsep=2pt,itemsep=1pt]
  \item Use eigenkets for measurement probabilities.
  \item Construct projectors $\hat{P}_a$ and verify $\hat{P}_a^2=\hat{P}_a$.
  \item Write spectral decomposition for discrete observables.
  \item State measurement postulate (Sakurai formulation).
\end{itemize}

\medskip\textbf{Tier differentiation:}
HS: ``filters'' and outcomes intuition --- Stern--Gerlach analogy.
Beg.\ UG: compute $P(a)=|\braket{a|\psi}|^2$ via projectors.
Adv.\ UG: degeneracy; resolution of identity in eigenbasis.
MSc: projection postulate subtleties; L\"uders rule.
PhD: spectral measures; PVMs and POVMs preview; measurement models.

\medskip\textbf{Assessment:}
Set~A --- compute probabilities and post-measurement states.
Set~B --- degeneracy, continuous spectrum cases.
Projects: (simple) spin-$\tfrac{1}{2}$ measurement simulation; (complex) L\"uders vs.\ von Neumann update comparison.
\end{tcolorbox}

\medskip
\begin{tcolorbox}[modulebox, title={Module~I.2 --- Lecture L05: Completeness, Change of Basis \& Matrix Mechanics}]
\textbf{90-min pacing:} (0--20~min) Resolution of identity; inserting $\hat{1}$;
(20--50~min) Basis expansion and change of representation;
(50--75~min) Matrix multiplication as inserted identities;
(75--90~min) Coordinate transforms; unitary basis changes.

\medskip\textbf{Core learning outcomes:}
\begin{itemize}[leftmargin=1.4em,topsep=2pt,itemsep=1pt]
  \item Use completeness $\sum_n\ket{n}\bra{n}=\hat{1}$ to evaluate any bracket.
  \item Transform between bases using unitary matrices.
  \item Derive matrix multiplication from operator composition.
  \item Express operator matrix elements $A_{mn}=\bra{m}\hat{A}\ket{n}$.
\end{itemize}

\medskip\textbf{Tier differentiation:}
HS: coordinate systems analogy, $2$D rotation.
Beg.\ UG: compute $A_{mn}$ and transform vectors between bases.
Adv.\ UG: show equivalence of abstract operator and matrix representation.
MSc: similarity vs.\ unitary equivalence; trace and spectrum invariants.
PhD: representation theory framing; direct integrals preview.

\medskip\textbf{Assessment:}
Set~A --- basis expansions, matrix elements.
Set~B --- change-of-basis proofs, unitary equivalence.
Projects: (simple) matrix mechanics calculator; (complex) equivalence of Schr\"odinger and Heisenberg representations.
\end{tcolorbox}

\medskip
\begin{tcolorbox}[modulebox, title={Module~I.2 --- Lecture L06: Continuous Bases: $\ket{x}$, $\ket{p}$ \& Delta Functions}]
\textbf{90-min pacing:} (0--20~min) Continuous completeness $\int\ket{x}\bra{x}dx=\hat{1}$;
(20--50~min) Wavefunction as $\psi(x)=\braket{x|\psi}$;
(50--75~min) Momentum basis; $\braket{x|p}=e^{ipx/\hbar}/\sqrt{2\pi\hbar}$;
(75--90~min) Fourier connection; delta function identities.

\medskip\textbf{Core learning outcomes:}
\begin{itemize}[leftmargin=1.4em,topsep=2pt,itemsep=1pt]
  \item Understand continuous-spectrum orthonormality $\braket{x|x'}=\delta(x-x')$.
  \item Derive wavefunction as projection onto position basis.
  \item Connect Fourier transform to $\braket{x|p}$ inner product.
  \item Handle delta function distributions carefully.
\end{itemize}

\medskip\textbf{Tier differentiation:}
HS: what ``continuum of outcomes'' means physically.
Beg.\ UG: basic integrals; qualitative delta function properties.
Adv.\ UG: careful delta manipulations; normalisation conventions.
MSc: distribution theory viewpoint; generalised eigenvectors.
PhD: rigged Hilbert spaces motivation; Gel'fand triple preview.

\medskip\textbf{Assessment:}
Set~A --- wavefunction integrals, Fourier transforms.
Set~B --- distribution theory; momentum-space wavefunctions.
Projects: (simple) Gaussian wavepacket in $x/p$ bases; (complex) prove Parseval--Plancherel from completeness.
\end{tcolorbox}

\medskip
\begin{tcolorbox}[modulebox, title={Module~I.2 --- Lecture L07: Measurement Theory: Postulates, Compatibility \& Uncertainty}]
\textbf{90-min pacing:} (0--20~min) Compatible observables and commutators $[\hat{A},\hat{B}]$;
(20--50~min) Uncertainty relation derivation (Robertson);
(50--75~min) Measurement update; expectation values;
(75--90~min) POVM preview; state tomography mention.

\medskip\textbf{Core learning outcomes:}
\begin{itemize}[leftmargin=1.4em,topsep=2pt,itemsep=1pt]
  \item Define compatible observables: $[\hat{A},\hat{B}]=0 \Leftrightarrow$ common eigenbasis.
  \item Derive Robertson--Schr\"odinger uncertainty inequality.
  \item Compute variances and uncertainty products.
  \item Distinguish uncertainty from measurement imprecision.
\end{itemize}

\medskip\textbf{Tier differentiation:}
HS: uncertainty as ``spread'' not instrument error.
Beg.\ UG: compute variances; simple commutators $[\hat{x},\hat{p}]=i\hbar$.
Adv.\ UG: derive Robertson--Schr\"odinger inequality from Cauchy--Schwarz.
MSc: POVMs, operational view; state tomography preview.
PhD: measurement models (von Neumann coupling), disturbance, contextuality.

\medskip\textbf{Assessment:}
Set~A --- commutators, variances, uncertainty products.
Set~B --- prove Robertson inequality; compatible observable theorems.
Projects: (simple) uncertainty product plotter; (complex) POVM vs.\ projective measurement comparison.
\end{tcolorbox}

\medskip
\begin{tcolorbox}[modulebox, title={Module~I.2 --- Lecture L08: Time Evolution \& Pictures of Quantum Mechanics}]
\textbf{90-min pacing:} (0--20~min) Unitary evolution $\hat{U}(t)=e^{-i\hat{H}t/\hbar}$;
(20--45~min) Schr\"odinger picture: states evolve;
(45--70~min) Heisenberg picture: operators evolve;
(70--90~min) Constants of motion; $[\hat{A},\hat{H}]=0\Leftrightarrow$ conserved.

\medskip\textbf{Core learning outcomes:}
\begin{itemize}[leftmargin=1.4em,topsep=2pt,itemsep=1pt]
  \item Derive $\hat{U}(t)$ from unitarity and time-translation invariance.
  \item Work fluently in both Schr\"odinger and Heisenberg pictures.
  \item Derive Heisenberg equation of motion $\dot{\hat{A}}_H=\tfrac{i}{\hbar}[\hat{H},\hat{A}_H]$.
  \item Identify constants of motion via commutator with $\hat{H}$.
\end{itemize}

\medskip\textbf{Tier differentiation:}
HS: exponential as ``repeated tiny evolution steps.''
Beg.\ UG: differentiate $\langle A\rangle$ with time; simple Hamiltonians.
Adv.\ UG: Heisenberg equation of motion; Ehrenfest's theorem.
MSc: time-ordering; interaction picture preview; formal solution.
PhD: Stone's theorem; one-parameter unitary groups; domain issues.

\medskip\textbf{Assessment:}
Set~A --- evolve states; compute $\langle A\rangle(t)$.
Set~B --- prove Heisenberg EOM; constants of motion theorems.
Projects: (simple) two-level Rabi oscillation; (complex) interaction picture and perturbation theory preview.
\end{tcolorbox}

\medskip
\begin{tcolorbox}[modulebox, title={Module~I.2 --- Lecture L09: Tensor Products: Composite Systems \& Entanglement}]
\textbf{90-min pacing:} (0--20~min) Tensor product space $\mathcal{H}_A\otimes\mathcal{H}_B$;
(20--50~min) Product vs.\ entangled states; operator action;
(50--75~min) Entanglement criterion; Schmidt decomposition;
(75--90~min) Partial measurement and reduced descriptions.

\medskip\textbf{Core learning outcomes:}
\begin{itemize}[leftmargin=1.4em,topsep=2pt,itemsep=1pt]
  \item Construct product states $\ket{i}_A\otimes\ket{j}_B$ and general bipartite states.
  \item Identify entanglement by failure of coefficient factorisation.
  \item Apply Schmidt decomposition to find Schmidt rank.
  \item Preview partial trace as reduced description of subsystem.
\end{itemize}

\medskip\textbf{Tier differentiation:}
HS: ``two systems = bigger space'' --- classical analogy and beyond.
Beg.\ UG: compute product basis components; tensor product of matrices.
Adv.\ UG: Schmidt decomposition; maximally entangled states.
MSc: partial trace introduced lightly; entanglement entropy preview.
PhD: entanglement measures, operational tasks, reading-based prompts.

\medskip\textbf{Assessment:}
Set~A --- product basis, operator action on $\mathcal{H}_A\otimes\mathcal{H}_B$.
Set~B --- Schmidt decomposition; Bell states analysis.
Projects: (simple) entanglement checker for 2-qubit states; (complex) Schmidt rank and entanglement entropy.
\end{tcolorbox}

\medskip
\begin{tcolorbox}[modulebox, title={Module~I.2 --- Lecture L10: Spin \& Finite-Dimensional Hilbert Spaces ($SU(2)$)}]
\textbf{90-min pacing:} (0--20~min) Spin-$\tfrac{1}{2}$ as 2D Hilbert space; $\ket{+},\ket{-}$ basis;
(20--50~min) Pauli matrices; spin operators $\sigma_x,\sigma_y,\sigma_z$;
(50--75~min) Rotation operators $e^{-i\theta\boldsymbol{\sigma}\cdot\hat{n}/2}$; Bloch sphere;
(75--90~min) Angular momentum commutation relations preview.

\medskip\textbf{Core learning outcomes:}
\begin{itemize}[leftmargin=1.4em,topsep=2pt,itemsep=1pt]
  \item Work with spin-$\tfrac{1}{2}$ as prototype finite-dimensional QM.
  \item Compute eigenkets and probabilities for all spin components.
  \item Apply rotation operators $U(\theta,\hat{n})=e^{-i\theta\boldsymbol{\sigma}\cdot\hat{n}/2}$.
  \item Connect $SU(2)$ group to physical rotations $SO(3)$.
\end{itemize}

\medskip\textbf{Tier differentiation:}
HS: two-state systems as ``qubits'' --- conceptual motivation.
Beg.\ UG: compute spin measurement probabilities; Stern--Gerlach results.
Adv.\ UG: rotation operators; verify $[\sigma_i,\sigma_j]=2i\varepsilon_{ijk}\sigma_k$.
MSc: representation theory touch; angular momentum algebra $\hat{J}_\pm$.
PhD: $SU(2)$ vs.\ $SO(3)$, projective representations, Wigner theorem link.

\medskip\textbf{Assessment:}
Set~A --- spin probabilities, expectation values.
Set~B --- rotation operators; angular momentum algebra.
Projects: (simple) Bloch sphere visualiser; (complex) $SU(2)\leftrightarrow SO(3)$ double-cover demonstration.
\end{tcolorbox}

\medskip
\begin{tcolorbox}[modulebox, title={Module~I.2 --- Lecture L11: Density Operators, Mixed States \& Partial Trace}]
\textbf{90-min pacing:} (0--20~min) Pure vs.\ mixed states;
$\hat{\rho}=\sum_i p_i\ket{\psi_i}\bra{\psi_i}$;
(20--50~min) $\langle\hat{A}\rangle=\operatorname{Tr}(\hat{\rho}\hat{A})$; properties of $\hat{\rho}$;
(50--75~min) Partial trace; reduced density matrices;
(75--90~min) Bloch sphere for spin-$\tfrac{1}{2}$; purity $\operatorname{Tr}(\hat{\rho}^2)$.

\medskip\textbf{Core learning outcomes:}
\begin{itemize}[leftmargin=1.4em,topsep=2pt,itemsep=1pt]
  \item Construct density operator for ensembles and reduced states.
  \item Compute expectation values via $\operatorname{Tr}(\hat{\rho}\hat{A})$.
  \item Apply partial trace to obtain subsystem state.
  \item Distinguish pure ($\operatorname{Tr}(\hat{\rho}^2)=1$) from mixed ($\operatorname{Tr}(\hat{\rho}^2)<1$) states.
\end{itemize}

\medskip\textbf{Tier differentiation:}
HS: ``uncertainty because we don't know'' vs.\ genuine quantum spread.
Beg.\ UG: compute trace in matrix form; verify $\hat{\rho}$ properties.
Adv.\ UG: pure vs.\ mixed; Bloch vector parameterisation of $2\times 2$ density matrices.
MSc: Kraus operators; CPTP maps preview; quantum channels.
PhD: open systems, master equations, quantum error correction preview.

\medskip\textbf{Assessment:}
Set~A --- density matrix construction, $\operatorname{Tr}(\hat{\rho}\hat{A})$.
Set~B --- partial trace, purity, Bloch representation.
Projects: (simple) density matrix visualiser for 2-level system; (complex) decoherence model via partial trace.
\end{tcolorbox}

\medskip
\begin{tcolorbox}[modulebox, title={Module~I.2 --- Lecture L12: Rigged Hilbert Spaces \& Mathematical Foundations}]
\textbf{90-min pacing:} (0--20~min) Why $\ket{x}$ is \emph{not} in the Hilbert space;
(20--50~min) Distributions and test functions; Schwartz space;
(50--75~min) Self-adjointness vs.\ Hermiticity revisited;
(75--90~min) What Dirac notation hides --- safe usage rules.

\medskip\textbf{Core learning outcomes:}
\begin{itemize}[leftmargin=1.4em,topsep=2pt,itemsep=1pt]
  \item Explain why $\braket{x|x'}=\delta(x-x')$ requires distribution theory.
  \item State the Gel'fand triple $\Phi\subset\mathcal{H}\subset\Phi'$.
  \item Distinguish self-adjoint from Hermitian (deficiency indices).
  \item Apply Dirac notation safely with awareness of its limits.
\end{itemize}

\medskip\textbf{Tier differentiation:}
HS: conceptual ``idealised states'' --- measuring infinitely precise position.
Beg.\ UG: cautionary rules for delta functions in calculations.
Adv.\ UG: examples of unbounded operators ($\hat{x},\hat{p}$); domain issues.
MSc: rigged Hilbert space triple $\Phi\subset\mathcal{H}\subset\Phi'$; nuclear spectral theorem.
PhD: spectral theorem in full measure-theoretic form; deficiency indices; self-adjoint extensions.

\medskip\textbf{Assessment:}
Set~A --- delta function identities; distribution manipulations.
Set~B --- self-adjoint extension problems; domain analysis.
Projects: (simple) ``regularisation'' demonstration; (complex) momentum operator on half-line --- self-adjoint extensions.
\end{tcolorbox}

%------------------------------------------------------------
\subsubsection*{Assessment Strategy}
\addcontentsline{toc}{subsubsection}{Module I.2 Assessment Strategy}
\label{sec:braket-assessment}

\begin{tcolorbox}[highlightbox]
\textbf{Per-Lecture Assessment Bundle} (every lecture):
concept check (10--15~min, 5--8 short conceptual questions);
Problem Set~A (simple, 6--10 computational exercises);
Problem Set~B (complex, 3--6 multi-step derivations and proofs);
mini-project simple (1--3 hours);
mini-project complex (5--10 hours, open-ended);
research questions well-defined (2--4 scoped prompts);
research questions fully open (1--2 frontier-style prompts).
\end{tcolorbox}

\medskip
\noindent\textbf{Course-level grading template:}
\begin{itemize}[leftmargin=1.5em]
  \item Weekly homework (best 10/12): 40\%
  \item Two mid-course syntheses (after L06 \& L10): 20\%
  \item Project (simple or complex track): 25\%
  \item Final oral/whiteboard exam or take-home: 15\%
\end{itemize}
Mixed cohorts: HS/UG $\to$ simple track; MSc/PhD $\to$ complex track.

%------------------------------------------------------------
\subsubsection*{Reading Strategy}
\addcontentsline{toc}{subsubsection}{Module I.2 Reading Strategy}

\noindent\textbf{Track A --- Core textbooks:}
Dirac, \textit{The Principles of Quantum Mechanics} (formalism, notation ethos);
Sakurai \& Napolitano, \textit{Modern Quantum Mechanics} (measurement, symmetry framing).

\noindent\textbf{Track B --- Support:}
Ballentine (postulates/measurement);
von Neumann (mathematical structure);
functional analysis references (PhD tier).

\noindent\textbf{Track C --- Historical sources:}
Born (1926) probability interpretation;
Dirac (1927) amplitude formalism;
von Neumann (1932) foundations;
Stone (1932) unitary groups;
Wigner, symmetry principles (MSc/PhD).

%------------------------------------------------------------
\subsubsection*{Expansion Plan --- Natural Extensions Beyond 12 Lectures}
\addcontentsline{toc}{subsubsection}{Module I.2 Expansion Plan}

\begin{center}
\begin{tabular}{lll}
\toprule
\textbf{Topic} & \textbf{Duration} & \textbf{Content}\\
\midrule
Angular Momentum \& Addition & 2--4 lectures & $\hat{J}_\pm$; Clebsch--Gordan; Wigner--Eckart theorem\\
Harmonic Oscillator \& Ladder Operators & 2--3 lectures & $\hat{a},\hat{a}^\dagger$; coherent states; squeezed states\\
Scattering Theory \& $S$-Matrix & 2--4 lectures & Lippmann--Schwinger; T-matrix; partial waves\\
Path Integrals & 2 lectures & Feynman formulation; propagator; connection to Series~II Module~II.2\\
Quantum Information Formalism & 2--4 lectures & Quantum channels; CPTP; quantum error correction\\
Advanced Measurement \& POVMs & 2 lectures & POVM elements; Naimark's theorem; informationally complete\\
\bottomrule
\end{tabular}
\end{center}

\newpage
%=============================================================
\subsection{Module I.3 --- Quantum Dynamics, Perturbation Theory, and Physical Systems}
\label{sec:mod-I3}
%=============================================================

\coursebadge{%
  \textbf{Scope:} Continues directly from Module~I.1 (\S2.4 onward). Covers the
  density matrix in full, bases, commutation relations, symmetries and Noether's
  theorem, the complete theory of time evolution (time-independent, time-dependent,
  and non-commuting), all three pictures of quantum mechanics, perturbation theory,
  open systems, and a survey of canonical example systems and physical effects.\\[2pt]
  \textbf{Duration:} 11 lectures $\times$ 90 minutes (Lecture~L01 continued + Lectures~L02--L11
  of the original Series~I programme).\\[2pt]
  \textbf{Prerequisite:} Module~I.1; Module~I.2 recommended as concurrent or prior study.
}

\subsubsection*{Module Overview}
\addcontentsline{toc}{subsubsection}{Module I.3 Overview}

\begin{center}
\begin{tabularx}{\textwidth}{clX}
\toprule
\textbf{Lec.} & \textbf{Title} & \textbf{Key Concepts}\\
\midrule
L01\textsuperscript{cont} & Pure/Mixed States and the Density Matrix (\S2.3--2.5) &
    Pure vs.\ mixed states, $\hat{\rho}$ properties, $\operatorname{Tr}\hat{\rho}=1$, purity, partial trace, entanglement entropy preview\\
L02 & Bases: Discrete and Continuous &
    Orthonormal bases, completeness relations, position/momentum representations, Dirac delta\\
L03 & Uncertainty and Commutation Relations &
    Robertson--Schr\"odinger inequality, compatible observables, CSCO, simultaneous eigenstates\\
L04 & Infinitesimal Generators, Symmetries, and Noether's Theorem &
    Lie algebras, Stone's theorem, Noether's theorem, discrete symmetries\\
L05 & Time Evolution I --- Time-Independent Hamiltonians &
    Time evolution operator $\hat{U}(t)$, T-symmetry, stationary states, Schr\"odinger equation\\
L06 & Time Evolution II --- Time-Dependent and Non-Commuting &
    Dyson series, time-ordered exponential $\mathcal{T}\exp$, Magnus expansion\\
L07 & Pictures of Quantum Mechanics and the Heisenberg Picture &
    Schr\"odinger/Heisenberg pictures, equation of motion, Ehrenfest, virial theorem\\
L08 & The Dirac (Interaction) Picture &
    Splitting $\hat{H}=\hat{H}_0+\hat{V}$, interaction picture, $S$-matrix\\
L09 & Perturbation Theory &
    Rayleigh--Schr\"odinger PT (degenerate and non-degenerate), Born--Oppenheimer, Feynman--Hellmann\\
L10 & Time-Dependent PT, Open Systems, and Many-Body Preview &
    Fermi's golden rule, Lindblad equation, decoherence, two-electron system\\
L11 & Example Systems and Physical Effects &
    Free particle, particle in a box, Bloch electron, harmonic oscillator,
    Rabi model, Zeeman/Stark/spin--orbit\\
\bottomrule
\end{tabularx}
\end{center}

\subsubsection*{Lecture-by-Lecture Syllabi}
\addcontentsline{toc}{subsubsection}{Module I.3 Lecture Syllabi}

\begin{tcolorbox}[modulebox, title={Module~I.3 --- Lecture L01 (cont.): Pure/Mixed States and the Density Matrix (\S2.3--2.5)}]
\textbf{90-min pacing:} (0--20~min) Motivation: why a single ket is insufficient to describe
sub-systems or statistical ensembles; pure state $\hat{\rho}=\ket{\psi}\bra{\psi}$ as
a special case;
(20--45~min) Mixed states: statistical mixture $\hat{\rho}=\sum_i p_i\ket{\psi_i}\bra{\psi_i}$;
physical distinction between classical ignorance and genuine quantum superposition;
(45--70~min) Properties of the density operator: $\hat{\rho}^\dagger=\hat{\rho}$,
$\operatorname{Tr}\hat{\rho}=1$, $\hat{\rho}\geq 0$; purity
$\gamma=\operatorname{Tr}\hat{\rho}^2\leq 1$ (equality iff pure);
expectation values $\langle\hat{A}\rangle=\operatorname{Tr}(\hat{\rho}\hat{A})$;
(70--90~min) Bipartite systems $\mathcal{H}_A\otimes\mathcal{H}_B$; partial trace
$\hat{\rho}_A=\operatorname{Tr}_B\hat{\rho}_{AB}$; entanglement entropy
$S=-\operatorname{Tr}(\hat{\rho}_A\ln\hat{\rho}_A)$ (preview).

\medskip\textbf{Core learning outcomes:}
\begin{itemize}[leftmargin=1.4em,topsep=2pt,itemsep=1pt]
  \item Construct $\hat{\rho}$ for pure states, mixtures, and subsystems.
  \item Verify the three defining properties ($\hat{\rho}^\dagger=\hat{\rho}$,
    $\operatorname{Tr}\hat{\rho}=1$, $\hat{\rho}\geq 0$) for given density matrices.
  \item Compute expectation values and purity; distinguish pure from mixed states.
  \item Apply partial trace to obtain reduced density matrices and compute entanglement entropy.
\end{itemize}

\medskip\textbf{Dual-track content:}
\begin{itemize}[leftmargin=1.4em,topsep=2pt,itemsep=1pt]
  \item \textit{Statistical/physical track:} Ensemble interpretation of $\hat{\rho}$;
    classical vs.\ quantum uncertainty; Bloch sphere picture for spin-$\tfrac{1}{2}$;
    why a mixed state cannot be ``purified'' within the same Hilbert space.
  \item \textit{Mathematical/algebraic track:} $\hat{\rho}$ as a positive semi-definite
    trace-class operator; spectral decomposition; partial trace as unique linear map
    satisfying $\operatorname{Tr}_B(\hat{\rho}_{AB}(\hat{1}\otimes\hat{O}_B))
    =\operatorname{Tr}(\hat{\rho}_B\hat{O}_B)$.
\end{itemize}

\medskip\textbf{Tier differentiation:}
HS: ``classical probability of being in a state'' vs.\ genuine quantum superposition ---
the key conceptual divide; coin-flip analogy and its limits.
Beg.\ UG: compute $\hat{\rho}$, $\operatorname{Tr}\hat{\rho}^2$, and $\langle\hat{A}\rangle$
for $2\times 2$ examples; verify properties numerically.
Adv.\ UG: prove $\operatorname{Tr}\hat{\rho}^2\leq 1$ from spectral decomposition;
Bloch vector parameterisation for spin-$\tfrac{1}{2}$; partial trace calculation.
MSc: Purification of mixed states; Schmidt decomposition connection;
entanglement entropy and its properties (sub-additivity preview).
PhD: Trace-class operators in infinite dimensions; von Neumann entropy;
quantum relative entropy; complete positivity preview.

\medskip\textbf{Assessment:}
Set~A --- construct $\hat{\rho}$ for explicit mixtures; compute purity and $\langle\hat{\sigma}_z\rangle$;
partial trace for 2-qubit Bell states.
Set~B --- prove $\operatorname{Tr}\hat{\rho}^2\leq(\operatorname{Tr}\hat{\rho})^2$ with equality
iff rank-1; show entanglement entropy vanishes for product states.
Projects: (simple) Bloch sphere trajectory under dephasing noise; (complex) purification of a given
mixed state and verification of entanglement entropy.
\end{tcolorbox}

\medskip
\begin{tcolorbox}[modulebox, title={Module~I.3 --- Lecture L02: Bases: Discrete and Continuous}]
\textbf{90-min pacing:} (0--20~min) Discrete orthonormal bases: $\braket{e_i|e_j}=\delta_{ij}$;
resolution of identity $\hat{1}=\sum_i\ket{e_i}\bra{e_i}$;
change of basis via unitary matrices; matrix representation in a given basis;
(20--45~min) Continuous bases: position eigenstates $\hat{x}\ket{x}=x\ket{x}$;
continuous orthonormality $\braket{x|x'}=\delta(x-x')$; completeness
$\hat{1}=\int dx\,\ket{x}\bra{x}$; wavefunction as $\psi(x)=\braket{x|\psi}$;
(45--70~min) Momentum representation: $\tilde{\psi}(p)=\braket{p|\psi}$;
$\braket{x|p}=e^{ipx/\hbar}/\sqrt{2\pi\hbar}$; Fourier transform as change of basis;
(70--90~min) The Dirac delta distribution: properties, Fourier representation,
regularisation; over-completeness and coherent states (preview).

\medskip\textbf{Core learning outcomes:}
\begin{itemize}[leftmargin=1.4em,topsep=2pt,itemsep=1pt]
  \item Use resolution of identity in discrete and continuous forms to evaluate any bracket.
  \item Derive the wavefunction $\psi(x)=\braket{x|\psi}$ as a projection.
  \item Perform basis changes between position and momentum representations via Fourier transforms.
  \item Manipulate delta functions and understand their distributional nature.
\end{itemize}

\medskip\textbf{Dual-track content:}
\begin{itemize}[leftmargin=1.4em,topsep=2pt,itemsep=1pt]
  \item \textit{Computational track:} Concrete calculation of Fourier transforms;
    Gaussian wavepackets in $x$ and $p$; Parseval--Plancherel theorem;
    practical use of $\hat{1}$ insertions to convert between representations.
  \item \textit{Mathematical track:} Generalised eigenvectors; distribution theory
    and Schwartz space; over-complete frames and the resolution of identity for
    coherent states; Gel'fand triple motivation.
\end{itemize}

\medskip\textbf{Tier differentiation:}
HS: ``snapshots'' of a state in different languages (position vs.\ momentum);
why the wavefunction $\psi(x)$ is just one way to represent the same ket.
Beg.\ UG: compute $\psi(x)$ from a given ket in a discrete basis; evaluate simple
Fourier transforms; use delta function properties.
Adv.\ UG: prove Parseval--Plancherel from completeness; handle delta function
carefully in operator matrix elements.
MSc: distribution theory viewpoint; generalised eigenvectors and their domain;
Fourier transform as unitary map on $L^2(\mathbb{R})$.
PhD: Gel'fand triple $\Phi\subset L^2\subset\Phi'$; nuclear spectral theorem;
coherent states and frame theory; wavelet transforms.

\medskip\textbf{Assessment:}
Set~A --- compute $\psi(x)$ and $\tilde{\psi}(p)$ for Gaussian and square states;
evaluate $\braket{x|\hat{p}|\psi}$ via completeness insertion.
Set~B --- prove Fourier transform is unitary on $L^2$; delta function identities and
distributional derivatives.
Projects: (simple) Gaussian wavepacket visualiser in $x$ and $p$; (complex) uncertainty
product $\Delta x\,\Delta p$ for a family of wavepackets and saturation condition.
\end{tcolorbox}

\begin{tcolorbox}[modulebox, title={Module~I.3 --- Lecture L03: Uncertainty and Commutation Relations}]
\textbf{90-min pacing:} (0--20~min) Commutators as the algebraic fingerprint of incompatibility;
$[\hat{x},\hat{p}]=i\hbar$; general commutator identities and Jacobi identity;
(20--45~min) The generalised uncertainty principle: variance definition
$(\Delta A)^2=\langle\hat{A}^2\rangle-\langle\hat{A}\rangle^2$;
Robertson inequality $\Delta A\,\Delta B\geq\tfrac{1}{2}|\langle[\hat{A},\hat{B}]\rangle|$
derived via Cauchy--Schwarz; Schr\"odinger's improvement with anticommutator term;
(45--65~min) Canonical commutation relations (CCRs):
$[\hat{x}_i,\hat{p}_j]=i\hbar\delta_{ij}$; Stone--von Neumann theorem (statement);
complete set of commuting observables (CSCO) and quantum numbers;
(65--90~min) Compatible and incompatible measurements; projection postulate;
sequential measurements; quantum Zeno effect; energy--time uncertainty
(careful treatment of $\Delta E\,\Delta t$).

\medskip\textbf{Core learning outcomes:}
\begin{itemize}[leftmargin=1.4em,topsep=2pt,itemsep=1pt]
  \item Derive Robertson's uncertainty inequality from the Cauchy--Schwarz inequality.
  \item Evaluate commutators for polynomial functions of $\hat{x}$ and $\hat{p}$.
  \item Define a CSCO and explain how it labels quantum states uniquely.
  \item Distinguish compatible from incompatible observables and their measurement implications.
\end{itemize}

\medskip\textbf{Dual-track content:}
\begin{itemize}[leftmargin=1.4em,topsep=2pt,itemsep=1pt]
  \item \textit{Physical track (Sakurai):} Uncertainty as an intrinsic property of
    quantum states, not of the measuring apparatus; Stern--Gerlach sequential
    measurement as the paradigm for incompatibility; energy--time uncertainty as
    a statement about state lifetime.
  \item \textit{Mathematical track:} Derivation of Robertson--Schr\"odinger inequality;
    Cauchy--Schwarz in Hilbert space; conditions for saturation (Gaussian states);
    Stone--von Neumann uniqueness of CCR representation.
\end{itemize}

\medskip\textbf{Tier differentiation:}
HS: ``you cannot know both position and momentum exactly'' --- why it is not
about clumsy measurements; wave analogy with bandwidth.
Beg.\ UG: compute $[\hat{x},\hat{p}^2]$, $[\hat{x}^2,\hat{p}]$ from CCR;
evaluate $\Delta x\,\Delta p$ for Gaussian.
Adv.\ UG: full derivation of Robertson inequality; conditions for saturation;
prove that $[\hat{A},\hat{B}]=0 \Leftrightarrow$ common eigenbasis (non-degenerate case).
MSc: Schr\"odinger's stronger uncertainty relation; CSCO for hydrogen;
Stone--von Neumann theorem and its physical content.
PhD: Weyl form of CCR; non-equivalent representations; quantum Zeno effect
rigorously; information-theoretic uncertainty relations (entropic).

\medskip\textbf{Assessment:}
Set~A --- evaluate commutators; compute variances for harmonic oscillator eigenstates;
verify Robertson inequality numerically.
Set~B --- derive Robertson inequality from scratch; prove saturation condition;
CSCO construction for a particle in 3D central potential.
Projects: (simple) uncertainty product plotter for parameterised family of states;
(complex) entropic uncertainty relation vs.\ Robertson for conjugate bases.
\end{tcolorbox}

\medskip
\begin{tcolorbox}[modulebox, title={Module~I.3 --- Lecture L04: Infinitesimal Generators, Symmetries, and Noether's Theorem}]
\textbf{90-min pacing:} (0--20~min) One-parameter unitary groups:
$\hat{U}(\theta)=e^{i\theta\hat{G}}$; Stone's theorem identifying the
generator $\hat{G}$; momentum as generator of spatial translations,
angular momentum of rotations, Hamiltonian of time translations;
(20--45~min) Symmetry of the Hamiltonian: $[\hat{H},\hat{U}]=0
\Leftrightarrow[\hat{H},\hat{G}]=0$; conservation of $\hat{G}$
in the Heisenberg picture as Noether's theorem;
examples: translation invariance $\to$ momentum conservation,
rotational invariance $\to$ angular momentum conservation;
(45--65~min) Discrete symmetries: parity $\hat{\Pi}$, charge conjugation $\hat{C}$,
time reversal $\hat{T}$ (Wigner anti-unitary); $CPT$ theorem (statement and significance);
(65--90~min) Symmetry breaking: explicit vs.\ spontaneous;
degenerate ground states; Goldstone modes (qualitative).

\medskip\textbf{Core learning outcomes:}
\begin{itemize}[leftmargin=1.4em,topsep=2pt,itemsep=1pt]
  \item Identify the generator of a given one-parameter unitary group.
  \item Prove Noether's theorem in quantum mechanics via the Heisenberg equation.
  \item Classify discrete symmetries and their representations (unitary vs.\ anti-unitary).
  \item Distinguish explicit from spontaneous symmetry breaking and describe physical consequences.
\end{itemize}

\medskip\textbf{Dual-track content:}
\begin{itemize}[leftmargin=1.4em,topsep=2pt,itemsep=1pt]
  \item \textit{Physical track:} Conservation laws as a direct consequence of symmetry;
    physical examples (translation $\to$ momentum, rotation $\to$ angular momentum,
    time $\to$ energy); spontaneous symmetry breaking in ferromagnets and superconductors (qualitative).
  \item \textit{Mathematical track:} Lie group -- Lie algebra correspondence via exponential map;
    Stone's theorem; adjoint action and structure constants;
    Wigner's theorem applied to discrete symmetries.
\end{itemize}

\medskip\textbf{Tier differentiation:}
HS: ``if physics looks the same after a transformation, something is conserved'' ---
Noether's theorem as the deepest connection between symmetry and conservation.
Beg.\ UG: verify $[\hat{H},\hat{p}]=0$ for translationally invariant $\hat{H}$;
check parity of simple wavefunctions.
Adv.\ UG: prove Noether's theorem from Heisenberg equation;
derive parity selection rules; show $\hat{\Pi}\hat{x}\hat{\Pi}^{-1}=-\hat{x}$.
MSc: Lie algebra generators; structure constants; $\mathfrak{su}(2)$ algebra;
full Wigner classification of discrete symmetries.
PhD: Projective representations; central extensions; spontaneous symmetry breaking
and Goldstone theorem; Wigner--Eckart theorem preview.

\medskip\textbf{Assessment:}
Set~A --- identify generators for rotations, translations, phase shifts;
verify conservation laws for given Hamiltonians; parity of wavefunctions.
Set~B --- prove Noether's theorem in the Heisenberg picture;
classify symmetries of the hydrogen Hamiltonian; analyse $CPT$ invariance.
Projects: (simple) symmetry classification tool for toy Hamiltonians;
(complex) spontaneous symmetry breaking in double-well potential ---
ground state degeneracy and perturbative lifting.
\end{tcolorbox}

\medskip
\begin{tcolorbox}[modulebox, title={Module~I.3 --- Lecture L05: Time Evolution I --- Time-Independent Hamiltonians}]
\textbf{90-min pacing:} (0--15~min) The time-evolution postulate:
$i\hbar\partial_t\ket{\psi(t)}=\hat{H}\ket{\psi(t)}$; linearity and
unitarity as fundamental requirements;
(15--40~min) The time evolution operator $\hat{U}(t,t_0)$: definition,
group property $\hat{U}(t_2,t_0)=\hat{U}(t_2,t_1)\hat{U}(t_1,t_0)$,
unitarity $\hat{U}^\dagger\hat{U}=\hat{1}$;
time-independent case $\hat{U}(t)=e^{-i\hat{H}t/\hbar}$ via spectral decomposition;
(40--65~min) Stationary states: $\hat{H}\ket{E_n}=E_n\ket{E_n}$;
general solution as superposition with phase factors $e^{-iE_n t/\hbar}$;
expectation values of time-independent operators in energy eigenstates;
(65--90~min) Time-reversal symmetry: anti-unitary $\hat{T}$;
Kramers degeneracy for half-integer spin;
$T$-invariant Hamiltonians; Schr\"odinger equation in position representation:
$\hat{p}=-i\hbar\nabla$; probability current $\mathbf{j}=\tfrac{\hbar}{2mi}(\psi^*\nabla\psi-\psi\nabla\psi^*)$;
continuity equation $\partial_t|\psi|^2+\nabla\cdot\mathbf{j}=0$.

\medskip\textbf{Core learning outcomes:}
\begin{itemize}[leftmargin=1.4em,topsep=2pt,itemsep=1pt]
  \item Derive $\hat{U}(t)=e^{-i\hat{H}t/\hbar}$ from the time-evolution postulate.
  \item Evolve an arbitrary initial state by expanding in the energy eigenbasis.
  \item Identify and exploit stationary state structure for computing expectation values.
  \item Derive the probability current and continuity equation in position representation.
\end{itemize}

\medskip\textbf{Dual-track content:}
\begin{itemize}[leftmargin=1.4em,topsep=2pt,itemsep=1pt]
  \item \textit{Wave mechanics track:} Schr\"odinger equation as a PDE; boundary conditions;
    probability current; stationary states as separable solutions $\psi(x,t)=\phi(x)e^{-iEt/\hbar}$.
  \item \textit{Abstract operator track:} $\hat{U}$ from Stone's theorem; spectral decomposition
    of $e^{-i\hat{H}t/\hbar}$; time translation group; anti-unitary time reversal
    and Kramers theorem.
\end{itemize}

\medskip\textbf{Tier differentiation:}
HS: ``the Hamiltonian is the quantum engine of time evolution'' --- energy eigenstates
oscillate simply; superpositions interfere as time goes on.
Beg.\ UG: evolve a two-level system; compute $\langle\hat{\sigma}_z\rangle(t)$;
derive Rabi oscillation frequency.
Adv.\ UG: full derivation of $\hat{U}$; probability current derivation;
prove unitarity of $e^{-i\hat{H}t/\hbar}$ for Hermitian $\hat{H}$.
MSc: time-reversal symmetry; Kramers degeneracy proof for half-integer spin;
compare time-independent and time-dependent treatments.
PhD: Stone's theorem; generator domain; self-adjoint extensions and their role in
determining the evolution; Zeno's paradox and its resolution.

\medskip\textbf{Assessment:}
Set~A --- evolve superposition states; compute $\langle\hat{E}\rangle(t)$,
$\langle\hat{x}\rangle(t)$; derive probability current for plane waves.
Set~B --- prove group property of $\hat{U}$; unitarity; derive continuity equation;
analyse Kramers degeneracy for spin-$\tfrac{1}{2}$ particle.
Projects: (simple) time evolution visualiser for a particle in a box;
(complex) Kramers degeneracy breaking by explicit $T$-violation and energy level
splitting analysis.
\end{tcolorbox}

\medskip
\begin{tcolorbox}[modulebox, title={Module~I.3 --- Lecture L06: Time Evolution II --- Time-Dependent and Non-Commuting Hamiltonians}]
\textbf{90-min pacing:} (0--20~min) Failure of $\hat{U}=e^{-i\hat{H}t/\hbar}$ when
$[\hat{H}(t_1),\hat{H}(t_2)]\neq 0$; integral equation for $\hat{U}(t,t_0)$;
iterative solution and the Dyson series;
(20--45~min) Time-ordered exponential
$\mathcal{T}\exp\!\bigl(-\tfrac{i}{\hbar}\int_{t_0}^{t}\hat{H}(t')dt'\bigr)$;
Dyson's time-ordering operator $\mathcal{T}$; first- and second-order terms explicitly;
(45--65~min) Piecewise-constant approximation: product formula
$\hat{U}\approx\prod_k e^{-i\hat{H}(t_k)\Delta t/\hbar}$;
connection to Trotter--Suzuki decomposition and quantum simulation;
(65--90~min) Magnus expansion: exponentiation of the Dyson series;
Magnus cumulants $\Omega_1,\Omega_2,\ldots$; unitarity preservation;
convergence radius; physical example: driven spin-$\tfrac{1}{2}$ in an
oscillating field; preview of Berry phase.

\medskip\textbf{Core learning outcomes:}
\begin{itemize}[leftmargin=1.4em,topsep=2pt,itemsep=1pt]
  \item Derive the Dyson series for a time-dependent Hamiltonian.
  \item Write the time-ordered exponential and evaluate its first two terms.
  \item Explain why the Magnus expansion preserves unitarity while the Dyson series does not term-by-term.
  \item Apply the piecewise-constant approximation to a simple driven system.
\end{itemize}

\medskip\textbf{Dual-track content:}
\begin{itemize}[leftmargin=1.4em,topsep=2pt,itemsep=1pt]
  \item \textit{Physical track:} Driven two-level system (Rabi problem exact vs.\ perturbative);
    rotating wave approximation origin; geometric phase accumulation as a physical effect
    of non-commutativity; quantum control applications.
  \item \textit{Mathematical track:} Time-ordering as a resummation of a Neumann series;
    Magnus expansion as the unique Hermitian exponent; Zassenhaus-type decompositions;
    convergence of Magnus expansion (Baker norm criterion).
\end{itemize}

\medskip\textbf{Tier differentiation:}
HS: ``when the Hamiltonian changes in time, the future depends on the exact history''
--- analogy with a spinning top being poked at different times.
Beg.\ UG: first-order Dyson term; transition probability formula;
numerical piecewise evolution for a two-level system.
Adv.\ UG: full Dyson series derivation; prove time-ordering is necessary;
compute $\Omega_1$ and $\Omega_2$ of Magnus expansion.
MSc: Magnus expansion convergence; Trotter--Suzuki product formula and its error;
connection to quantum simulation circuits.
PhD: Magnus expansion in Lie algebraic language; Floquet theory for periodic
Hamiltonians; geometric phases and holonomy in the Magnus framework.

\medskip\textbf{Assessment:}
Set~A --- compute first-order Dyson term for oscillating Hamiltonian;
piecewise evolution of a qubit; transition probability to first order.
Set~B --- derive Dyson series from integral equation; prove $\Omega_1$ of
Magnus expansion; convergence criterion application.
Projects: (simple) numerical Dyson vs.\ Magnus evolution comparison for driven qubit;
(complex) Floquet Hamiltonian for periodically driven two-level system and
quasi-energy spectrum.
\end{tcolorbox}

\medskip
\begin{tcolorbox}[modulebox, title={Module~I.3 --- Lecture L07: Pictures of Quantum Mechanics and the Heisenberg Picture}]
\textbf{90-min pacing:} (0--15~min) The equivalence-class of pictures:
physical predictions (transition probabilities) are picture-independent;
the Schr\"odinger picture reviewed --- states evolve, operators fixed;
(15--40~min) The Heisenberg picture: operators carry the time dependence,
$\hat{A}_H(t)=\hat{U}^\dagger(t)\hat{A}\hat{U}(t)$; states static
$\ket{\psi_H}=\ket{\psi(t_0)}$; Heisenberg equation of motion
$\dot{\hat{A}}_H = \tfrac{i}{\hbar}[\hat{H},\hat{A}_H]+(\partial_t\hat{A})_H$;
(40--65~min) Consequences: Ehrenfest's theorem
$\tfrac{d}{dt}\langle\hat{x}\rangle=\langle\hat{p}\rangle/m$,
$\tfrac{d}{dt}\langle\hat{p}\rangle=-\langle\nabla V\rangle$;
virial theorem; constants of motion $[\hat{H},\hat{A}]=0\Rightarrow\dot{\hat{A}}_H=0$;
connection to Noether (Lecture~L04);
(65--90~min) Canonical quantisation in the Heisenberg picture:
classical Poisson brackets $\to$ commutators,
$\{f,g\}_{\rm PB}\to\tfrac{1}{i\hbar}[\hat{f},\hat{g}]$;
Dirac's quantisation rule; ordering ambiguities.

\medskip\textbf{Core learning outcomes:}
\begin{itemize}[leftmargin=1.4em,topsep=2pt,itemsep=1pt]
  \item Transform operators and states between Schr\"odinger and Heisenberg pictures.
  \item Derive and apply the Heisenberg equation of motion.
  \item Prove and use Ehrenfest's theorem and the virial theorem.
  \item Apply Dirac's canonical quantisation rule and identify operator ordering issues.
\end{itemize}

\medskip\textbf{Dual-track content:}
\begin{itemize}[leftmargin=1.4em,topsep=2pt,itemsep=1pt]
  \item \textit{Physical track (Sakurai):} Heisenberg picture as the quantum analogue
    of classical mechanics --- operators play the role of dynamical variables;
    Ehrenfest's theorem as the quantum--classical correspondence;
    constants of motion as quantum-mechanical conservation laws.
  \item \textit{Mathematical/structural track:} Heisenberg equation as a derivation on the
    operator algebra; relationship to Lie bracket structure; canonical quantisation as
    a map from Poisson manifold to non-commutative algebra.
\end{itemize}

\medskip\textbf{Tier differentiation:}
HS: ``in one picture the states move, in the other the measuring devices move'' ---
the physics is the same; analogy with active vs.\ passive coordinate transformations.
Beg.\ UG: transform $\hat{x}(t)$ and $\hat{p}(t)$ for a harmonic oscillator;
verify Ehrenfest's theorem numerically.
Adv.\ UG: full derivation of Heisenberg EOM; virial theorem proof;
solve $\hat{x}_H(t)$ and $\hat{p}_H(t)$ for particle in a quadratic potential.
MSc: Dirac quantisation rule; operator ordering ambiguities (Weyl ordering);
connection to classical integrability and action--angle variables.
PhD: Algebraic quantum mechanics; Heisenberg picture in $C^*$-algebra language;
Moyal bracket and Wigner function as semi-classical bridge.

\medskip\textbf{Assessment:}
Set~A --- solve Heisenberg EOM for harmonic oscillator; compute $\langle\hat{x}\rangle(t)$
and $\langle\hat{p}\rangle(t)$; verify Ehrenfest for a cubic potential.
Set~B --- prove virial theorem; apply canonical quantisation to a classical system;
analyse ordering ambiguities for $\hat{x}^2\hat{p}$.
Projects: (simple) comparison of Schr\"odinger and Heisenberg pictures for a
driven harmonic oscillator; (complex) Moyal bracket vs.\ quantum commutator for
polynomial observables.
\end{tcolorbox}

\medskip
\begin{tcolorbox}[modulebox, title={Module~I.3 --- Lecture L08: The Dirac (Interaction) Picture}]
\textbf{90-min pacing:} (0--15~min) Motivation: perturbative treatment of $\hat{V}(t)$
while exploiting exact knowledge of $\hat{H}_0$ dynamics;
splitting $\hat{H}(t)=\hat{H}_0+\hat{V}(t)$;
(15--40~min) Definition of the interaction picture:
states $\ket{\psi_I(t)}=e^{i\hat{H}_0 t/\hbar}\ket{\psi_S(t)}$;
operators $\hat{V}_I(t)=e^{i\hat{H}_0 t/\hbar}\hat{V}\,e^{-i\hat{H}_0 t/\hbar}$;
interaction picture Schr\"odinger equation $i\hbar\partial_t\ket{\psi_I}=\hat{V}_I(t)\ket{\psi_I}$;
(40--65~min) Interaction picture propagator $\hat{U}_I(t,t_0)$;
Dyson series in the interaction picture; perturbative expansion of transition amplitudes;
$S$-matrix $=\hat{U}_I(\infty,-\infty)$; physical interpretation as scattering operator;
(65--90~min) Comparison table of all three pictures: state and operator evolution,
equations of motion, preferred use-cases; bridge to time-dependent perturbation theory (L09--L10).

\medskip\textbf{Core learning outcomes:}
\begin{itemize}[leftmargin=1.4em,topsep=2pt,itemsep=1pt]
  \item Transform states and operators into and out of the interaction picture.
  \item Derive the interaction picture Schr\"odinger equation from the full equation.
  \item Write the Dyson series in the interaction picture to any desired order.
  \item Connect the $S$-matrix to the interaction picture propagator.
\end{itemize}

\medskip\textbf{Dual-track content:}
\begin{itemize}[leftmargin=1.4em,topsep=2pt,itemsep=1pt]
  \item \textit{Physical track:} Interaction picture as the natural language of
    perturbation theory and scattering; $S$-matrix elements as transition amplitudes;
    relationship to measurable cross-sections and decay rates.
  \item \textit{Mathematical track:} Interaction picture as a unitary equivalence
    transformation; two-sided action of free evolution; composition with Dyson series
    to produce $S$-matrix; time-ordering and causality.
\end{itemize}

\medskip\textbf{Tier differentiation:}
HS: ``park the easy part of the Hamiltonian in the operators, and let the hard part
drive the states'' --- interaction picture as a practical bookkeeping device.
Beg.\ UG: transform a simple interaction $\hat{V}=\lambda\hat{x}$ into the
interaction picture for a harmonic oscillator $\hat{H}_0$; first-order Dyson term.
Adv.\ UG: full transformation procedure; Dyson series to second order;
prove equivalence of all three pictures for physical predictions.
MSc: $S$-matrix and cross-sections; optical theorem;
Lippmann--Schwinger equation preview; Feynman diagram structure.
PhD: $S$-matrix in QFT context; adiabatic switching; Gell-Mann--Low theorem;
LSZ reduction formula (conceptual).

\medskip\textbf{Assessment:}
Set~A --- transform given operators and states to interaction picture;
compute first-order transition amplitude for a harmonic perturbation.
Set~B --- derive interaction picture EOM; prove unitarity of $S$-matrix;
second-order Dyson series for a constant perturbation.
Projects: (simple) three-picture comparison for a driven two-level system;
(complex) $S$-matrix for two-level system under short pulse: exact vs.\ first/second-order.
\end{tcolorbox}

\medskip
\begin{tcolorbox}[modulebox, title={Module~I.3 --- Lecture L09: Perturbation Theory}]
\textbf{90-min pacing:} (0--20~min) Setup: $\hat{H}=\hat{H}_0+\lambda\hat{V}$;
known eigenstates $\hat{H}_0\ket{n^{(0)}}=E_n^{(0)}\ket{n^{(0)}}$;
perturbative ansatz $E_n=\sum_k\lambda^k E_n^{(k)}$,
$\ket{n}=\sum_k\lambda^k\ket{n^{(k)}}$; non-degenerate case;
(20--45~min) Rayleigh--Schr\"odinger PT: first-order energy $E_n^{(1)}=\bra{n^{(0)}}\hat{V}\ket{n^{(0)}}$;
first-order state correction $\ket{n^{(1)}}=\sum_{m\neq n}\tfrac{\bra{m^{(0)}}\hat{V}\ket{n^{(0)}}}{E_n^{(0)}-E_m^{(0)}}\ket{m^{(0)}}$;
second-order energy; convergence and breakdown conditions;
(45--65~min) Degenerate perturbation theory: identify degenerate subspace;
block diagonalise $\hat{V}$ within it; secular equation; correct zeroth-order states;
worked example (hydrogen Stark effect);
(65--90~min) Born--Oppenheimer separation;
Feynman--Hellmann theorem $\partial_\lambda E_n=\langle n|\partial_\lambda\hat{H}|n\rangle$;
adiabatic theorem (statement); Ritz variational principle as complementary tool.

\medskip\textbf{Core learning outcomes:}
\begin{itemize}[leftmargin=1.4em,topsep=2pt,itemsep=1pt]
  \item Derive first- and second-order energy corrections in non-degenerate RSPT.
  \item Compute first-order state corrections and apply them to matrix elements.
  \item Handle degenerate perturbation theory by diagonalising $\hat{V}$ in the degenerate subspace.
  \item Apply the Feynman--Hellmann theorem and the Ritz variational principle.
\end{itemize}

\medskip\textbf{Dual-track content:}
\begin{itemize}[leftmargin=1.4em,topsep=2pt,itemsep=1pt]
  \item \textit{Physical/applied track:} Hydrogen atom in electric field (Stark effect);
    fine structure as perturbation; Born--Oppenheimer in molecular physics;
    variational quantum eigensolvers (preview).
  \item \textit{Mathematical/formal track:} Kato--Rellich theorem on analytic perturbations;
    convergence of Rayleigh--Schr\"odinger series; resolvent operator
    $\hat{R}_0(z)=(z-\hat{H}_0)^{-1}$ and Brillouin--Wigner PT.
\end{itemize}

\medskip\textbf{Tier differentiation:}
HS: ``a small push changes energy levels slightly'' --- first-order PT as the
quantum analogue of Taylor expansion for energy levels.
Beg.\ UG: first-order energy corrections for simple perturbations;
linear Stark effect for hydrogen $n=2$ by diagonalisation.
Adv.\ UG: second-order energy; convergence conditions; full degenerate PT
procedure; Feynman--Hellmann theorem proof.
MSc: Kato--Rellich theorem; Brillouin--Wigner PT; Born--Oppenheimer rigorously;
adiabatic theorem statement and proof sketch.
PhD: Resolvent formalism; Dyson equation for the resolvent;
divergence of perturbation series and resummation; PT for resonances.

\medskip\textbf{Assessment:}
Set~A --- first-order corrections for anharmonic oscillator $\lambda\hat{x}^4$;
linear and quadratic Stark effect for hydrogen; variational bound for ground state.
Set~B --- second-order energy for harmonic oscillator shifted by linear term;
degenerate PT for $n=2$ hydrogen in electric field; Feynman--Hellmann for HO.
Projects: (simple) PT vs.\ exact energy levels for anharmonic oscillator as a function
of $\lambda$; (complex) Brillouin--Wigner vs.\ Rayleigh--Schr\"odinger comparison and
convergence.
\end{tcolorbox}

\medskip
\begin{tcolorbox}[modulebox, title={Module~I.3 --- Lecture L10: Time-Dependent PT, Open Systems, and Many-Body Preview}]
\textbf{90-min pacing:} (0--20~min) Time-dependent perturbation theory:
first-order transition amplitude
$c_f^{(1)}(t)=-\tfrac{i}{\hbar}\int_{t_0}^{t}dt'\,e^{i\omega_{fi}t'}\bra{f}\hat{V}(t')\ket{i}$;
transition probability $|c_f^{(1)}|^2$; sudden and adiabatic limits;
(20--40~min) Fermi's golden rule:
continuous spectrum; density of states $\rho(E_f)$;
$\Gamma_{i\to f}=\tfrac{2\pi}{\hbar}|\bra{f}\hat{V}\ket{i}|^2\rho(E_f)$;
derivation from first-order TDPT via sinc function limit; selection rules;
(40--65~min) Open quantum systems: closed (unitary) vs.\ open (environment coupling);
Lindblad master equation $\dot{\hat{\rho}}=-\tfrac{i}{\hbar}[\hat{H},\hat{\rho}]
+\sum_k({\hat{L}_k\hat{\rho}\hat{L}_k^\dagger-\tfrac{1}{2}\{\hat{L}_k^\dagger\hat{L}_k,\hat{\rho}\}})$;
jump operators; decay channels; dephasing;
(65--90~min) Decoherence and the measurement problem:
environment-induced superselection; pointer states;
many-body preview: two-electron system, exchange interaction, singlet and triplet states,
para- and ortho-helium.

\medskip\textbf{Core learning outcomes:}
\begin{itemize}[leftmargin=1.4em,topsep=2pt,itemsep=1pt]
  \item Derive Fermi's golden rule from first-order TDPT.
  \item Write the Lindblad master equation and identify physical contributions of each term.
  \item Explain decoherence as environment-induced suppression of off-diagonal density matrix elements.
  \item Construct two-electron spin states (singlet and triplet) and explain their exchange origin.
\end{itemize}

\medskip\textbf{Dual-track content:}
\begin{itemize}[leftmargin=1.4em,topsep=2pt,itemsep=1pt]
  \item \textit{Physical/applied track:} Fermi's golden rule in atomic spectroscopy,
    nuclear decay, and solid-state (electron--phonon scattering); Lindblad equation in
    quantum optics (spontaneous emission) and quantum computing (gate errors);
    decoherence as the bridge between quantum and classical.
  \item \textit{Mathematical/structural track:} Derivation of Fermi's golden rule rigorously;
    Lindblad equation from complete positivity and the GKSL theorem;
    pointer states as preferred bases selected by the environment.
\end{itemize}

\medskip\textbf{Tier differentiation:}
HS: ``open systems'' as a leaky box --- energy and information escaping to the environment;
radioactive decay as Fermi's golden rule in action.
Beg.\ UG: compute $\Gamma_{i\to f}$ for a constant perturbation;
simple Lindblad evolution (amplitude damping) for a qubit.
Adv.\ UG: derive Fermi's golden rule from TDPT; verify Lindblad equation
preserves $\operatorname{Tr}\hat{\rho}=1$ and positivity; singlet/triplet construction.
MSc: GKSL theorem and complete positivity; Kraus representation;
decoherence in the spin-boson model; Born--Markov approximation.
PhD: Nakajima--Zwanzig projection technique; non-Markovian dynamics;
Floquet--Lindblad equation; entanglement in many-body systems.

\medskip\textbf{Assessment:}
Set~A --- compute $\Gamma_{i\to f}$ for an atom in an oscillating field;
simple Lindblad evolution; singlet/triplet state construction and symmetry.
Set~B --- derive Fermi's golden rule; prove Lindblad equation trace-preserving;
exchange interaction from antisymmetry of two-electron wavefunction.
Projects: (simple) decoherence visualiser: off-diagonal decay under Lindblad dephasing;
(complex) Born--Markov derivation of Lindblad equation for a qubit coupled to a harmonic bath.
\end{tcolorbox}

\medskip
\begin{tcolorbox}[modulebox, title={Module~I.3 --- Lecture L11: Example Systems and Physical Effects}]
\textbf{90-min pacing:} (0--20~min) Free particle and particle in a box:
plane waves and group velocity; wave packets; infinite square well: quantisation,
zero-point energy, parity; simple Bloch electron: periodic potential, Bloch's theorem,
crystal momentum, band structure (elementary), Brillouin zone;
(20--40~min) Harmonic oscillator: ladder operators $\hat{a},\hat{a}^\dagger$;
spectrum $E_n=(n+\tfrac{1}{2})\hbar\omega$; coherent states; Fock states;
driven and damped anharmonic oscillator: resonance, anharmonic perturbations
$\lambda\hat{x}^3,\lambda\hat{x}^4$;
(40--60~min) Rabi model and two-level systems:
Jaynes--Cummings Hamiltonian; Rabi oscillations; rotating wave approximation;
$\pi$- and $\pi/2$-pulses; minimal coupling Hamiltonian
$\hat{H}=\tfrac{1}{2m}(\hat{\mathbf{p}}-e\mathbf{A})^2+e\phi$;
gauge invariance; Landau levels;
(60--90~min) Zeeman and Stark effects: normal and anomalous Zeeman effect;
linear and quadratic Stark effect; selection rules;
spin--orbit coupling: relativistic origin, Thomas precession,
$\hat{H}_{\rm SO}=\xi(r)\hat{\mathbf{L}}\cdot\hat{\mathbf{S}}$;
$j$-$j$ vs.\ $L$-$S$ coupling; fine structure of hydrogen.

\medskip\textbf{Core learning outcomes:}
\begin{itemize}[leftmargin=1.4em,topsep=2pt,itemsep=1pt]
  \item Solve the harmonic oscillator algebraically using ladder operators and derive the spectrum.
  \item Derive Rabi oscillations and apply the rotating wave approximation.
  \item Compute Zeeman and Stark energy shifts using first-order perturbation theory.
  \item Explain spin--orbit coupling physically and compute fine-structure splittings.
\end{itemize}

\medskip\textbf{Dual-track content:}
\begin{itemize}[leftmargin=1.4em,topsep=2pt,itemsep=1pt]
  \item \textit{Physical survey track:} Each system as a worked example of the formalism
    developed in L01--L10; emphasis on physical interpretation (zero-point energy, coherent
    states as classical-like states, Rabi cycling as quantum control, Zeeman spectroscopy);
    connections to atomic and condensed matter experiments.
  \item \textit{Mathematical/algebraic track:} Ladder operator algebra and its generalisation;
    Bloch's theorem from translation symmetry; minimal coupling and gauge invariance;
    angular momentum addition and Clebsch--Gordan coefficients in spin--orbit coupling.
\end{itemize}

\medskip\textbf{Tier differentiation:}
HS: ``quantum mechanics in action'' --- particle in a box as the simplest real system;
why atoms have sharp spectral lines (energy quantisation); laser pulses as $\pi$-pulses.
Beg.\ UG: solve harmonic oscillator with $\hat{a},\hat{a}^\dagger$; compute
first-order Zeeman splitting; Rabi oscillation frequency.
Adv.\ UG: derive Bloch's theorem from translation invariance;
rotating wave approximation derivation; anomalous Zeeman from spin;
spin--orbit splitting in hydrogen via first-order PT.
MSc: coherent states as eigenstates of $\hat{a}$; over-completeness; Landau levels;
$j$-$j$ vs.\ $L$-$S$ coupling comparison; Thomas precession derivation.
PhD: Jaynes--Cummings model exactly; dressed states; polaritons;
gauge invariance and Aharonov--Bohm effect; Dirac equation origin of spin--orbit.

\medskip\textbf{Assessment:}
Set~A --- harmonic oscillator matrix elements $\bra{m}\hat{x}^2\ket{n}$;
Rabi oscillation period; linear Stark shift of hydrogen ground state.
Set~B --- derive Bloch's theorem; anharmonic oscillator perturbation to second order;
spin--orbit fine structure of hydrogen $2p$ level.
Projects: (simple) energy level diagram for hydrogen with Zeeman, Stark, and spin--orbit;
(complex) Jaynes--Cummings model dynamics: collapse and revival of Rabi oscillations.
\end{tcolorbox}

%------------------------------------------------------------
\subsubsection*{Assessment Strategy}
\addcontentsline{toc}{subsubsection}{Module I.3 Assessment Strategy}

\begin{tcolorbox}[highlightbox]
\textbf{Per-Lecture Assessment Bundle} (every lecture):
concept check (10--15~min, 5--8 short conceptual questions);
Problem Set~A (simple, 6--10 computational exercises);
Problem Set~B (complex, 3--6 multi-step derivations and proofs);
mini-project simple (1--3 hours);
mini-project complex (5--10 hours, open-ended);
research questions well-defined (2--4 scoped prompts);
research questions fully open (1--2 frontier-style prompts).
\end{tcolorbox}

\medskip
\noindent\textbf{Course-level grading template:}
\begin{itemize}[leftmargin=1.5em]
  \item Weekly homework (best 9/11): 40\%
  \item Mid-course synthesis (after L06): 20\%
  \item Project (simple or complex track): 25\%
  \item Final oral/whiteboard exam or take-home: 15\%
\end{itemize}
Mixed cohorts: HS/UG $\to$ simple track; MSc/PhD $\to$ complex track.

%------------------------------------------------------------
\subsubsection*{Reading Strategy}
\addcontentsline{toc}{subsubsection}{Module I.3 Reading Strategy}

\noindent\textbf{Track A --- Core textbooks:}
Sakurai \& Napolitano, \textit{Modern Quantum Mechanics} Ch.~2--5 (dynamics, symmetry, perturbation theory);
Shankar, \textit{Principles of Quantum Mechanics} (perturbation theory and example systems).

\noindent\textbf{Track B --- Support:}
Messiah, \textit{Quantum Mechanics} (density matrix, open systems);
Cohen-Tannoudji, Diu \& Lalo\"e (time-dependent PT, Fermi's golden rule).

\noindent\textbf{Track C --- Specialised:}
Lindblad (1976) original paper on master equations (MSc/PhD);
Born \& Oppenheimer (1927) original paper;
Bloch (1929) for Bloch's theorem in solids.

\newpage

%#############################################################
%  SERIES II SEPARATOR PAGE
%#############################################################
\newpage
\thispagestyle{empty}
\pagecolor{darkblue}
\color{white}
\vspace*{2.2cm}
\begin{center}
  {\color{accentgold}\rule{11cm}{1.5pt}}\\[0.5cm]
  {\fontsize{11}{14}\selectfont\bfseries\color{accentgold}\MakeUppercase{Series II}}\\[0.3cm]
  {\fontsize{26}{32}\selectfont\bfseries Advanced Quantum Mechanics}\\[0.4cm]
  {\color{accentgold}\rule{11cm}{1.5pt}}\\[1.0cm]
  {\large\itshape Seven modules --- 84 lectures $\times$ 90 minutes}\\[0.5cm]
  {\normalsize\itshape Audience: Graduate students}\\[0.3cm]
  {\normalsize\itshape Prerequisite: Series~I or equivalent}\\[1.4cm]

  \begin{tabular}{@{}clr@{}}
    \toprule
    \textbf{Module} & \textbf{Title} & \textbf{Lectures}\\
    \midrule
    \textbf{II.1} & Ground-State Approximation Methods & 12\\[3pt]
                  & \multicolumn{2}{@{}l}{\small\color{accentgold}%
                    Ritz; linear variational; Hartree--Fock; coupled-cluster;}\\
                  & \multicolumn{2}{@{}l}{\small\color{accentgold}%
                    QMC (variational \& diffusion); DMRG; tensor networks; natural orbitals.}\\[6pt]
    \textbf{II.2} & Path Integrals in Quantum Mechanics & 12\\[3pt]
                  & \multicolumn{2}{@{}l}{\small\color{accentgold}%
                    Feynman postulate; Gaussian integrals; WKB; imaginary-time formalism;}\\
                  & \multicolumn{2}{@{}l}{\small\color{accentgold}%
                    coherent states; instantons; gauge fields; Berry phase.}\\[6pt]
    \textbf{II.3} & Theory of Angular Momentum & 12\\[3pt]
                  & \multicolumn{2}{@{}l}{\small\color{accentgold}%
                    $\mathfrak{su}(2)$ algebra; orbital and spin AM; Clebsch--Gordan;}\\
                  & \multicolumn{2}{@{}l}{\small\color{accentgold}%
                    Wigner--Eckart theorem; hydrogen; fine/hyperfine structure; $D$-matrices.}\\[6pt]
    \textbf{II.4} & Second Quantisation and Introductory Many-Body Theory & 12\\[3pt]
                  & \multicolumn{2}{@{}l}{\small\color{accentgold}%
                    Fock space; normal ordering; Wick's theorem; HF in 2nd quantisation;}\\
                  & \multicolumn{2}{@{}l}{\small\color{accentgold}%
                    Hubbard model; Green's functions; Dyson equation; BCS theory.}\\[6pt]
    \textbf{II.5} & Non-Equilibrium Quantum Many-Body Green's Functions & 12\\[3pt]
                  & \multicolumn{2}{@{}l}{\small\color{accentgold}%
                    Keldysh contour; NEGF; Kadanoff--Baym equations; self-energies;}\\
                  & \multicolumn{2}{@{}l}{\small\color{accentgold}%
                    quantum transport; Meir--Wingreen formula; TDDFT connection.}\\[6pt]
    \textbf{II.6} & Light--Matter Interaction, Linear Response, and Spectroscopy & 12\\[3pt]
                  & \multicolumn{2}{@{}l}{\small\color{accentgold}%
                    Minimal coupling; electric dipole approximation; Kubo formula;}\\
                  & \multicolumn{2}{@{}l}{\small\color{accentgold}%
                    linear response; dielectric function; optical spectroscopy; Casimir effect.}\\[6pt]
    \textbf{II.7} & Quantum Statistical Mechanics in Equilibrium & 12\\[3pt]
                  & \multicolumn{2}{@{}l}{\small\color{accentgold}%
                    Density matrix in stat.\ mech.; grand canonical; Bose--Einstein condensation;}\\
                  & \multicolumn{2}{@{}l}{\small\color{accentgold}%
                    Fermi liquid; Matsubara formalism; phase transitions; quantum criticality.}\\[4pt]
    \midrule
    & \textbf{Total} & \textbf{84}\\
    \bottomrule
  \end{tabular}

  \vfill
  {\color{accentgold}\rule{11cm}{0.8pt}}\\[0.4cm]
  {\small\itshape\color{white!80!darkblue}%
    Each module: 12 lectures $\times$ 90 minutes $\cdot$
    Prerequisite: Series~I $\cdot$ Research-oriented depth}\\[0.3cm]
  {\color{accentgold}\rule{11cm}{0.8pt}}
\end{center}
\newpage
\nopagecolor
\color{black}

%#############################################################
\section{Series II --- Advanced Quantum Mechanics}
\label{sec:seriesII}
%#############################################################

\begin{tcolorbox}[seriesbox, title={Series~II Overview}]
\textbf{Prerequisite:} Completion of Series~I or equivalent.\\
\textbf{Duration:} Each module comprises 12 lectures $\times$ 90 minutes.\\
\textbf{Core aim:} Systematically deepen the formalism and connect quantum
mechanics to modern research areas.\\[4pt]
\begin{tabular}{llr}
\toprule
\textbf{Module} & \textbf{Title} & \textbf{Lectures}\\
\midrule
II.1 & Ground-State Approximation Methods & 12\\
II.2 & Path Integrals in Quantum Mechanics & 12\\
II.3 & Theory of Angular Momentum & 12\\
II.4 & Second Quantisation and Introductory Many-Body Theory & 12\\
II.5 & Non-Equilibrium Quantum Many-Body Green's Functions & 12\\
II.6 & Light--Matter Interaction, Linear Response, and Spectroscopy & 12\\
II.7 & Quantum Statistical Mechanics in Equilibrium & 12\\
\midrule
 & \textbf{Total} & \textbf{84}\\
\bottomrule
\end{tabular}
\end{tcolorbox}

%=============================================================
\subsection{Module II.1 --- Ground-State Approximation Methods}
\label{sec:mod-II1}
%=============================================================

%------------------------------------------------------------
\subsubsection*{Module Overview}
\addcontentsline{toc}{subsubsection}{Module II.1 Overview}

\begin{center}
\begin{tabularx}{\textwidth}{clX}
\toprule
\textbf{Lec.} & \textbf{Title} & \textbf{Key Concepts}\\
\midrule
L01 & Variational Principle & Ritz method; Rayleigh quotient; upper bound theorem; trial states\\
L02 & Linear Variational Method & Secular determinant; basis expansion; H\"uckel MO theory\\
L03 & Hartree and Hartree--Fock & Mean-field; Fock operator; self-consistency; Koopman's theorem\\
L04 & Coupled-Cluster Methods & Cluster operator $e^{\hat{T}}$; CCSD(T); connected diagrams; size consistency\\
L05 & Quantum Monte Carlo I & Variational QMC; importance sampling; Metropolis algorithm\\
L06 & Quantum Monte Carlo II & Diffusion QMC; fixed-node approximation; sign problem\\
L07 & Density Matrix Renormalisation Group & Entanglement area laws; SVD; renormalisation group idea\\
L08 & Tensor Networks & MPS, MERA; contraction algorithms; bond dimension\\
L09 & Natural Orbitals & 1-RDM; NO expansion; occupation numbers; multi-configurational methods\\
L10 & Multi-Reference Methods & CASSCF; MRCI; NEVPT2; active space selection\\
L11 & Benchmark Systems & H$_2$, LiH, N$_2$ dissociation; accuracy hierarchy; error analysis\\
L12 & Method Selection in Practice & Cost scaling; applicability regimes; crossover criteria\\
\bottomrule
\end{tabularx}
\end{center}

\subsubsection*{Five-Tier Layered Pedagogy}
\addcontentsline{toc}{subsubsection}{Module II.1 Five-Tier Pedagogy}

\begin{center}
\begin{tabular}{llp{9cm}}
\toprule
\textbf{Tier} & \textbf{Cohort} & \textbf{Characteristic content}\\
\midrule
Tier~1 & HS       & Conceptual motivation; energy minimisation intuition; analogy with curve fitting\\
Tier~2 & Beg.\ UG & Definitions; matrix secular equations; basic HF; worked 2-electron examples\\
Tier~3 & Adv.\ UG & Rigorous derivations; convergence proofs; algorithm implementation; error analysis\\
Tier~4 & MSc      & Advanced methods (CC, QMC, DMRG); diagrammatic techniques; active space design\\
Tier~5 & PhD      & Frontiers: stochastic CC, FCIQMC, tensor hypercontraction, cost-scaling optimisation\\
\bottomrule
\end{tabular}
\end{center}

\subsubsection*{Lecture-by-Lecture Syllabi}
\addcontentsline{toc}{subsubsection}{Module II.1 Lecture Syllabi}

% ------ L01 -----------------------------------------------
\begin{tcolorbox}[modulebox, title={Module~II.1 --- Lecture L01: The Variational Principle}]
\textbf{90-min pacing:} (0--20~min) Motivation: exact diagonalisation is exponentially hard;
variational methods as the first systematic route around this wall;
statement of the variational theorem: for any normalised trial state $\ket{\phi}$,
$\langle\phi|\hat{H}|\phi\rangle \geq E_0$, with equality iff $\ket{\phi}=\ket{0}$;
(20--45~min) Proof from spectral decomposition; Rayleigh quotient
$R[\phi]=\langle\phi|\hat{H}|\phi\rangle/\langle\phi|\phi\rangle$;
optimisation as a constrained minimisation; Lagrange multiplier and the stationarity
condition $\delta R=0 \Leftrightarrow \hat{H}\ket{\phi}=E\ket{\phi}$;
(45--70~min) Excited states: Hylleraas--Undheim--MacDonald theorem
(eigenvalues of a truncated Hamiltonian interlace the exact eigenvalues);
upper bounds on excited-state energies;
Temple's lower bound;
(70--90~min) Practical trial states: Gaussian-type orbitals;
Pad\'e approximants; neural-network wavefunctions (VMC preview);
error analysis: how a variational error $\delta E_{\rm var}$ in the energy
corresponds to a first-order error in the wavefunction.

\medskip\textbf{Core learning outcomes:}
\begin{itemize}[leftmargin=1.4em,topsep=2pt,itemsep=1pt]
  \item State and prove the variational theorem from spectral decomposition.
  \item Compute the Rayleigh quotient for a trial wavefunction and minimise it analytically.
  \item State the Hylleraas--Undheim theorem and apply it to excited-state bounds.
  \item Explain why a second-order error in the wavefunction gives only a fourth-order energy error.
\end{itemize}

\medskip\textbf{Dual-track content:}
\begin{itemize}[leftmargin=1.4em,topsep=2pt,itemsep=1pt]
  \item \textit{Physical/intuitive track:} Energy minimisation as ``best guess in a restricted
    class of states''; analogy with curve-fitting subject to a norm constraint; Gaussian
    orbitals as computationally convenient trial states.
  \item \textit{Mathematical track:} Ritz--Galerkin theory; min-max characterisation of eigenvalues;
    Temple's lower bound; error propagation from wavefunction to energy.
\end{itemize}

\medskip\textbf{Tier differentiation:}
HS: ``the ground-state energy is the lowest possible expectation value of the Hamiltonian'' ---
variational principle as quantum minimisation.
Beg.\ UG: compute $\langle\hat{H}\rangle$ for Gaussian trial state in a harmonic potential;
optimise over the width parameter $\alpha$.
Adv.\ UG: full proof of variational theorem; Hylleraas--Undheim theorem; Temple lower bound.
MSc: error analysis; Pad\'e and Hylleraas-type trial functions; multi-parameter optimisation
with gradient and Hessian.
PhD: neural-network wavefunctions; stochastic reconfiguration; natural gradient descent;
connection to variational Monte Carlo.

\medskip\textbf{Assessment:}
Set~A --- Rayleigh quotient for Gaussian trial in hydrogen; optimise $\alpha$; bound on $E_0$.
Set~B --- prove variational theorem; derive stationarity condition; Hylleraas--Undheim theorem proof.
Projects: (simple) variational parameter scan for helium with $Z_{\rm eff}$;
(complex) multi-parameter Hylleraas wavefunction for helium and convergence to exact energy.
\end{tcolorbox}

\medskip

% ------ L02 -----------------------------------------------
\begin{tcolorbox}[modulebox, title={Module~II.1 --- Lecture L02: The Linear Variational Method}]
\textbf{90-min pacing:} (0--20~min) Linear ansatz:
$\ket{\phi}=\sum_{k=1}^N c_k\ket{\chi_k}$ for a fixed basis $\{\ket{\chi_k}\}$;
Rayleigh quotient in matrix form;
Hamiltonian matrix $H_{kl}=\braket{\chi_k|\hat{H}|\chi_l}$ and
overlap matrix $S_{kl}=\braket{\chi_k|\chi_l}$;
(20--45~min) Secular (generalised eigenvalue) equation $\mathbf{H}\mathbf{c}=E\mathbf{S}\mathbf{c}$;
Cholesky decomposition of $\mathbf{S}$ to convert to standard eigenvalue problem;
$N$ variational eigenvalues as upper bounds to exact eigenvalues (Hylleraas--Undheim);
orthonormalisation of basis via L\"owdin symmetric or canonical procedures;
(45--65~min) H\"uckel molecular orbital theory as a paradigm application:
$\pi$ system of ethylene, butadiene, benzene;
H\"uckel parameters $\alpha$, $\beta$; secular determinant; MO energies and coefficients;
aromaticity from H\"uckel rule $4n+2$;
(65--90~min) Extended H\"uckel and limitations; basis set convergence in general;
non-orthogonal basis subtleties; condition number of $\mathbf{S}$; linear dependence and
near-linear dependence; canonical truncation.

\medskip\textbf{Core learning outcomes:}
\begin{itemize}[leftmargin=1.4em,topsep=2pt,itemsep=1pt]
  \item Set up and solve the secular equation for a given basis.
  \item Apply H\"uckel MO theory to simple $\pi$ systems and interpret results.
  \item Orthonormalise a non-orthogonal basis via L\"owdin or canonical procedures.
  \item Diagnose and handle linear dependence in a basis set.
\end{itemize}

\medskip\textbf{Dual-track content:}
\begin{itemize}[leftmargin=1.4em,topsep=2pt,itemsep=1pt]
  \item \textit{Chemical/applied track:} H\"uckel MO energies and their relation to
    reactivity; aromatic stabilisation energy; frontier MO theory (HOMO/LUMO); extended H\"uckel.
  \item \textit{Mathematical track:} Generalised eigenvalue problem; Cholesky and L\"owdin
    orthogonalisation; condition number; linear dependence detection via SVD of overlap matrix.
\end{itemize}

\medskip\textbf{Tier differentiation:}
HS: H\"uckel MO theory for benzene --- why $\pi$ electrons delocalise and stabilise the ring.
Beg.\ UG: set up and solve $2\times 2$ secular equation; H\"uckel for ethylene and allyl.
Adv.\ UG: Hylleraas--Undheim theorem for variational eigenvalues; L\"owdin orthogonalisation derivation.
MSc: condition number analysis; canonical truncation; extended H\"uckel parameterisation.
PhD: connection to finite-element Galerkin methods; adaptively contracted basis sets;
DIIS convergence acceleration.

\medskip\textbf{Assessment:}
Set~A --- secular equation for particle-in-a-box with polynomial basis;
H\"uckel MO diagram for butadiene and benzene.
Set~B --- prove Hylleraas--Undheim theorem for linear variational method;
L\"owdin orthogonalisation derivation; condition number of Gaussian overlap matrix.
Projects: (simple) H\"uckel MO solver for arbitrary planar $\pi$ systems;
(complex) L\"owdin orthogonalisation vs.\ canonical truncation: numerical comparison
for near-linearly-dependent Gaussian basis.
\end{tcolorbox}

\medskip

% ------ L03 -----------------------------------------------
\begin{tcolorbox}[modulebox, title={Module~II.1 --- Lecture L03: Hartree and Hartree--Fock Theory}]
\textbf{90-min pacing:} (0--20~min) The $N$-electron problem; exponential wall;
Hartree product ansatz $\Psi=\prod_i\phi_i(\mathbf{r}_i)$; variational optimisation
gives Hartree equations $[\hat{h}+\hat{v}^{\rm H}_i]\phi_i=\varepsilon_i\phi_i$;
Hartree potential as classical mean field; self-consistency loop;
(20--50~min) Antisymmetry requirement: Slater determinant;
Hartree--Fock energy from variational principle on Slater determinant;
Fock operator $\hat{F}_i=\hat{h}_i+\hat{J}_i-\hat{K}_i$;
Coulomb $\hat{J}$ and exchange $\hat{K}$ operators defined and compared;
HF equations $\hat{F}\phi_i=\varepsilon_i\phi_i$;
canonical vs.\ localised orbitals;
(50--70~min) Self-consistent field (SCF) algorithm: Roothaan--Hall equations
$\mathbf{F}\mathbf{C}=\mathbf{S}\mathbf{C}\boldsymbol{\varepsilon}$ in an AO basis;
density matrix $\mathbf{P}=2\mathbf{C}_{\rm occ}\mathbf{C}_{\rm occ}^\dagger$;
convergence and DIIS; Koopman's theorem ($-\varepsilon_i \approx$ ionisation energy);
(70--90~min) Properties of HF: size consistency; variational; missing correlation;
correlation energy $E_{\rm corr}=E_{\rm exact}-E_{\rm HF}$;
restricted (RHF), unrestricted (UHF), and restricted-open-shell (ROHF);
spin contamination in UHF.

\medskip\textbf{Core learning outcomes:}
\begin{itemize}[leftmargin=1.4em,topsep=2pt,itemsep=1pt]
  \item Derive the Hartree and Hartree--Fock equations from a variational principle.
  \item Define the Coulomb and exchange operators and state their physical meaning.
  \item Describe the SCF algorithm and the Roothaan--Hall equations.
  \item Apply Koopman's theorem and state the limitations of HF theory.
\end{itemize}

\medskip\textbf{Dual-track content:}
\begin{itemize}[leftmargin=1.4em,topsep=2pt,itemsep=1pt]
  \item \textit{Physical/chemical track:} Mean-field picture; exchange as a purely quantum
    effect from antisymmetry; HF as the starting point for all post-HF methods;
    correlation hole; Fermi hole vs.\ Coulomb hole.
  \item \textit{Mathematical/formal track:} Functional derivative of the HF energy;
    Brillouin's theorem; Koopman's theorem derivation; relation of HF to DFT
    (Kohn--Sham as an orbital-optimised density-functional analogue).
\end{itemize}

\medskip\textbf{Tier differentiation:}
HS: ``each electron moves in the average field of all others'' --- mean field as an
approximation to avoid solving a coupled 3$N$-dimensional problem exactly.
Beg.\ UG: Hartree equations for helium; two-electron integrals; SCF iteration by hand.
Adv.\ UG: HF equations from variational principle; Brillouin's theorem; Koopman's theorem proof.
MSc: Roothaan--Hall in AO basis; DIIS convergence; symmetry breaking in UHF; RHF vs.\ UHF for H$_2$.
PhD: Brueckner orbitals; natural orbitals from 1-RDM; HF as saddle point of full CI energy;
orbital-optimised MP2.

\medskip\textbf{Assessment:}
Set~A --- HF energy of helium by direct integration; Koopman's theorem for first ionisation;
identify Coulomb vs.\ exchange terms in a given two-electron integral.
Set~B --- derive HF equations from $\delta E_{\rm HF}=0$; prove Brillouin's theorem;
show SCF density matrix idempotency $\mathbf{P}^2=\mathbf{P}$ at convergence.
Projects: (simple) HF SCF loop for $\rm H_2$ in a minimal (STO-3G) basis with convergence plot;
(complex) RHF vs.\ UHF dissociation curve for $\rm H_2$: spin contamination and correlation energy.
\end{tcolorbox}

\medskip

% ------ L04 -----------------------------------------------
\begin{tcolorbox}[modulebox, title={Module~II.1 --- Lecture L04: Coupled-Cluster Methods}]
\textbf{90-min pacing:} (0--20~min) Limitations of truncated CI: size inconsistency;
Brillouin--Wigner vs.\ Rayleigh--Schr\"odinger CI; size consistency defined;
the exponential ansatz: $\ket{\Psi_{\rm CC}}=e^{\hat{T}}\ket{\Phi_0}$ with cluster
operator $\hat{T}=\hat{T}_1+\hat{T}_2+\cdots$;
why the exponential structure guarantees size consistency;
(20--50~min) CC amplitude equations: project $\bra{\Phi_i^a}e^{-\hat{T}}\hat{H}e^{\hat{T}}\ket{\Phi_0}=0$
(doubles) and similarly for singles; Baker--Campbell--Hausdorff expansion of $e^{-\hat{T}}\hat{H}e^{\hat{T}}$
terminates at quartic in $\hat{T}$ (for two-body $\hat{H}$);
CCSD: $\hat{T}=\hat{T}_1+\hat{T}_2$; $\mathcal{O}(N^6)$ scaling;
(50--70~min) Perturbative triples: CCSD(T) --- the ``gold standard'';
$\mathcal{O}(N^7)$ scaling; relative cost and accuracy;
equation-of-motion CC (EOM-CC) for excited states; CCSD(T)-F12 for basis set
convergence acceleration;
(70--90~min) Diagrammatic representation of CC equations: Goldstone diagrams;
connected vs.\ disconnected; Wick's theorem applied to CC;
local correlation methods: PNO-CCSD(T), DLPNO-CCSD(T) and linear scaling.

\medskip\textbf{Core learning outcomes:}
\begin{itemize}[leftmargin=1.4em,topsep=2pt,itemsep=1pt]
  \item Explain why the exponential ansatz ensures size consistency.
  \item Derive the CCSD amplitude equations via BCH expansion.
  \item State the computational scaling of CCSD and CCSD(T) and place them in the accuracy hierarchy.
  \item Describe PNO-based local correlation and its scaling advantage.
\end{itemize}

\medskip\textbf{Dual-track content:}
\begin{itemize}[leftmargin=1.4em,topsep=2pt,itemsep=1pt]
  \item \textit{Computational chemistry track:} CCSD(T) as the de facto gold standard;
    benchmark accuracy for thermochemistry (sub-kJ/mol); EOM-CC for spectra;
    DLPNO-CCSD(T) for large molecules; software packages (ORCA, MOLPRO, CFOUR).
  \item \textit{Many-body theory track:} Diagrammatic CC; linked-cluster theorem;
    CC as a resummation of MBPT; connection to coupled-cluster Green's functions (CCGF).
\end{itemize}

\medskip\textbf{Tier differentiation:}
HS: ``instead of writing the wavefunction as a sum, write it as an exponential'' ---
why this small change guarantees that two well-separated molecules have independent energies.
Beg.\ UG: CCD amplitude equations in schematic form; size consistency demonstration for He$+$He.
Adv.\ UG: BCH expansion to quartic; derive CCSD equations; perturbative triples derivation.
MSc: EOM-CCSD; CCSD(T)-F12; PNO-CCSD(T) and local correlation ideas.
PhD: Diagrammatic CC; CCGF and ionisation potentials; transcorrelated CC; FCIQMC-tailored CC.

\medskip\textbf{Assessment:}
Set~A --- verify size consistency of CCSD for He dimer; CCD energy in schematic notation;
rank CCSD(T) in the accuracy hierarchy.
Set~B --- derive BCH expansion of $e^{-\hat{T}}\hat{H}e^{\hat{T}}$ to quartic; linked-cluster theorem.
Projects: (simple) CCSD vs.\ FCI energy curve for H$_2$ as function of bond length;
(complex) EOM-CCSD excitation energies for ethylene vs.\ TDDFT and experiment.
\end{tcolorbox}

\medskip

% ------ L05 -----------------------------------------------
\begin{tcolorbox}[modulebox, title={Module~II.1 --- Lecture L05: Quantum Monte Carlo I --- Variational QMC}]
\textbf{90-min pacing:} (0--20~min) Monte Carlo integration: basic idea; central limit theorem;
error $\propto 1/\sqrt{N_{\rm samples}}$; importance sampling to reduce variance;
application to high-dimensional integrals (the QMC motivation);
(20--45~min) Variational QMC (VMC): evaluate $E_{\rm var}=\langle\Psi_T|\hat{H}|\Psi_T\rangle/
\langle\Psi_T|\Psi_T\rangle$ stochastically; local energy $E_L(\mathbf{R})=\hat{H}\Psi_T/\Psi_T$;
variance minimisation vs.\ energy minimisation;
the Metropolis--Hastings algorithm: proposal, acceptance ratio
$A=\min(1,|\Psi_T(\mathbf{R}')|^2/|\Psi_T(\mathbf{R})|^2)$;
detailed balance and ergodicity;
(45--70~min) Trial wavefunction design: Slater--Jastrow form
$\Psi_T=\mathcal{A}[\phi]\exp(J)$;
Jastrow factor: electron--electron, electron--nuclear, three-body terms;
cusp conditions: $\partial_r E_L$ finite requires $(\partial_r\Psi_T/\Psi_T)|_{r=0}$ to
cancel the $1/r$ divergence ($s$-wave cusp: $a_{ee}=1/2$ for unlike spin, $1/4$ for like);
(70--90~min) Parameter optimisation: stochastic reconfiguration (SR) and the natural
gradient; ADAM and gradient descent in stochastic setting; correlated sampling;
VMC scaling: $\mathcal{O}(N^3)$--$\mathcal{O}(N^4)$ for determinant evaluation.

\medskip\textbf{Core learning outcomes:}
\begin{itemize}[leftmargin=1.4em,topsep=2pt,itemsep=1pt]
  \item Explain importance sampling and the Metropolis algorithm.
  \item Compute the local energy for a Slater--Jastrow trial wavefunction.
  \item Derive and apply the electron--electron cusp condition.
  \item Compare energy minimisation and variance minimisation for parameter optimisation.
\end{itemize}

\medskip\textbf{Dual-track content:}
\begin{itemize}[leftmargin=1.4em,topsep=2pt,itemsep=1pt]
  \item \textit{Computational/algorithmic track:} Metropolis algorithm implementation;
    autocorrelation time and statistical efficiency; blocking analysis for error estimation;
    GPU acceleration; practical VMC workflows.
  \item \textit{Mathematical/stochastic track:} Markov chain theory; detailed balance;
    ergodicity; central limit theorem for correlated samples; stochastic reconfiguration
    as natural gradient descent on the manifold of quantum states.
\end{itemize}

\medskip\textbf{Tier differentiation:}
HS: ``roll the dice millions of times to estimate an integral in 3$N$ dimensions'' ---
Monte Carlo as a solution to the curse of dimensionality.
Beg.\ UG: Metropolis walk for a harmonic oscillator; blocking analysis; variance estimate.
Adv.\ UG: Slater--Jastrow VMC; cusp conditions derivation; variance minimisation.
MSc: stochastic reconfiguration; natural gradient; correlated sampling for forces.
PhD: neural-network Jastrow and backflow; transcorrelated VMC; ab initio QMC for solids.

\medskip\textbf{Assessment:}
Set~A --- Metropolis walk for $|\psi_0|^2$ in HO; local energy for Gaussian trial; cusp check.
Set~B --- derive electron cusp condition; prove detailed balance for Metropolis; variance minimisation
stationarity condition.
Projects: (simple) VMC for helium with one-parameter Jastrow; convergence with $N_{\rm samples}$;
(complex) stochastic reconfiguration optimisation of multi-parameter Jastrow for LiH.
\end{tcolorbox}

\medskip

% ------ L06 -----------------------------------------------
\begin{tcolorbox}[modulebox, title={Module~II.1 --- Lecture L06: Quantum Monte Carlo II --- Diffusion QMC}]
\textbf{90-min pacing:} (0--20~min) From VMC to exact: the diffusion QMC (DMC) idea;
imaginary-time Schr\"odinger equation $-\partial_\tau\Psi=(\hat{H}-E_T)\Psi$;
large-$\tau$ projection onto ground state; analogy with diffusion + branching process;
Green's function $G(\mathbf{R}',\mathbf{R};\tau)$ as the propagator;
(20--45~min) DMC algorithm: walkers (configurations) undergo diffusion steps $\Delta\tau$;
branching (replication/death) with weight $\exp[-(E_L(\mathbf{R})-E_T)\Delta\tau]$;
importance sampling with trial function: modified Green's function;
drift velocity $\mathbf{F}=\nabla\ln|\Psi_T|$; guided diffusion;
(45--65~min) The fermion sign problem: positive and negative walkers annihilate;
fixed-node approximation: constrain walker crossings of nodal surface of $\Psi_T$;
fixed-node energy as an upper bound; quality of nodes determines accuracy;
release-node (unphysical) and its failure for electrons;
(65--90~min) Pseudopotentials in DMC: non-locality and T-moves; relativistic effects;
lattice regularised DMC; twist-averaged boundary conditions for periodic systems;
scaling: $\mathcal{O}(N^3)$; DMC for solids --- computational workflow and accuracy benchmarks.

\medskip\textbf{Core learning outcomes:}
\begin{itemize}[leftmargin=1.4em,topsep=2pt,itemsep=1pt]
  \item Derive the imaginary-time Schr\"odinger equation and explain ground-state projection.
  \item Describe the DMC algorithm with importance sampling and branching.
  \item Explain the sign problem and the fixed-node approximation.
  \item Assess the accuracy of DMC for atoms and molecules relative to CCSD(T).
\end{itemize}

\medskip\textbf{Dual-track content:}
\begin{itemize}[leftmargin=1.4em,topsep=2pt,itemsep=1pt]
  \item \textit{Physical/computational track:} DMC as the most accurate practical
    method for large systems; solid-state applications (transition metal oxides, ice polymorphs);
    comparison with CCSD(T) for benchmark sets; practical workflow with CASINO or QMCPack.
  \item \textit{Mathematical track:} Feynman--Kac formula; Green's function Monte Carlo;
    variational fixed-node theorem; node optimisation; relation to path-integral MC.
\end{itemize}

\medskip\textbf{Tier differentiation:}
HS: ``walkers explore space like diffusing particles, and the ones in low-energy regions
survive'' --- projection onto the ground state as a survival-of-the-fittest process.
Beg.\ UG: one-dimensional DMC for a particle in a box; branching and population control.
Adv.\ UG: fixed-node theorem proof; importance sampling derivation; timestep bias analysis.
MSc: T-moves for non-local pseudopotentials; twist-averaged boundary conditions; node optimisation.
PhD: Feynman--Kac formula; lattice regularisation; backflow nodes; auxiliary field QMC.

\medskip\textbf{Assessment:}
Set~A --- 1D DMC for HO; branching weight calculation; timestep extrapolation.
Set~B --- prove fixed-node variational theorem; derive importance-sampled Green's function.
Projects: (simple) DMC for helium: fixed-node energy vs.\ VMC and exact;
(complex) node quality study --- HF vs.\ CASSCF nodes for N$_2$ at various bond lengths.
\end{tcolorbox}

\medskip

% ------ L07 -----------------------------------------------
\begin{tcolorbox}[modulebox, title={Module~II.1 --- Lecture L07: Density Matrix Renormalisation Group}]
\textbf{90-min pacing:} (0--20~min) Exponential wall revisited: why the Hilbert space
dimension is $d^N$ for a spin-$\tfrac{1}{2}$ chain; entanglement as the key to efficient
representation; area law vs.\ volume law for ground states of gapped local Hamiltonians;
entanglement entropy $S = -\operatorname{Tr}(\hat{\rho}_A \ln\hat{\rho}_A)$
and its area law $S \leq c \cdot L^{d-1}$ in $d$ spatial dimensions;
(20--45~min) Schmidt decomposition and its cost: $\ket{\Psi}=\sum_\alpha \lambda_\alpha
\ket{\alpha}_L\ket{\alpha}_R$; singular value decomposition (SVD);
truncation to bond dimension $m$: approximation error bounded by discarded weights
$\epsilon = \sum_{\alpha>m}\lambda_\alpha^2$;
the renormalisation group idea: keep the $m$ most relevant states;
(45--70~min) Matrix Product State (MPS) representation:
$\ket{\Psi}=\sum_{\{s_i\}}\text{Tr}[A^{s_1}A^{s_2}\cdots A^{s_N}]\ket{s_1\cdots s_N}$;
bond dimension $m$ as the key parameter; DMRG as variational optimisation over MPS;
two-site DMRG algorithm: superblock Hamiltonian; Lanczos diagonalisation; SVD truncation;
(70--90~min) DMRG in quantum chemistry (DMRG-SCF / DMRG-CASSCF):
orbital ordering and entanglement; mutual information $I_{ij}$;
active space beyond CAS(18,18); scaling $\mathcal{O}(m^3 N + m^2 N d^2)$;
convergence with $m$ and comparison with CASSCF/MRCI.

\medskip\textbf{Core learning outcomes:}
\begin{itemize}[leftmargin=1.4em,topsep=2pt,itemsep=1pt]
  \item State the area law and explain why it justifies MPS representations.
  \item Perform a Schmidt decomposition and estimate the SVD truncation error.
  \item Describe the MPS ansatz and the two-site DMRG algorithm.
  \item Explain DMRG-CASSCF and its advantage over conventional CASSCF.
\end{itemize}

\medskip\textbf{Dual-track content:}
\begin{itemize}[leftmargin=1.4em,topsep=2pt,itemsep=1pt]
  \item \textit{Quantum chemistry/applications track:} DMRG as a large-active-space method;
    orbital selection via mutual information; applications to transition metal complexes,
    $f$-electron systems, and spin chains; software: ITensor, CheMPS2, Block2.
  \item \textit{Quantum information/mathematical track:} MPS as a variational class;
    area-law proof sketch (Hastings); entanglement spectrum; connection to topological order.
\end{itemize}

\medskip\textbf{Tier differentiation:}
HS: ``instead of keeping all $2^N$ amplitudes, keep only the most entangled ones'' ---
DMRG as smart truncation guided by quantum information.
Beg.\ UG: Schmidt decomposition for a 4-site chain; SVD truncation and error.
Adv.\ UG: MPS canonical forms; two-site DMRG algorithm; area law statement.
MSc: DMRG-CASSCF; orbital ordering optimisation; mutual information; convergence with $m$.
PhD: DMRG as tensor network; time-evolving block decimation (TEBD); finite-temperature DMRG;
DMRG for periodic systems.

\medskip\textbf{Assessment:}
Set~A --- Schmidt decomposition for a 4-qubit state; truncation error estimate;
MPS representation of a product state.
Set~B --- prove area-law scaling for 1D gapped systems (Hastings, sketch);
derive the variational condition in two-site DMRG.
Projects: (simple) DMRG energy vs.\ bond dimension $m$ for Heisenberg chain;
(complex) DMRG-CASSCF active space selection via mutual information for Cr$_2$ dimer.
\end{tcolorbox}

\medskip

% ------ L08 -----------------------------------------------
\begin{tcolorbox}[modulebox, title={Module~II.1 --- Lecture L08: Tensor Networks}]
\textbf{90-min pacing:} (0--20~min) MPS as a special case of a general tensor network (TN);
graphical notation: tensors as nodes, contracted indices as edges;
review of MPS canonical forms (left/right canonical);
why generalising beyond 1D requires new network geometries;
(20--45~min) Tree tensor networks (TTN): branching structure; efficient contraction;
multi-scale entanglement renormalisation ansatz (MERA): disentanglers and isometries;
causal cone; critical systems and logarithmic entanglement;
Projected Entangled Pair States (PEPS) for 2D systems: exponentially hard contraction;
practical approximations (boundary MPS, Monte Carlo sampling of PEPS);
(45--65~min) Contraction algorithms: optimal contraction order ($\mathcal{NP}$-hard in
general); DMRG as MPS contraction; tensor train decomposition;
Tucker decomposition and its connection to HOSVD;
tensor hypercontraction (THC) in quantum chemistry: $\mathcal{O}(N^4)\to\mathcal{O}(N^2)$
for two-electron integrals;
(65--90~min) Tensor network methods in machine learning: tensor trains as efficient
neural network layers; Born machine; quantum-inspired ML;
connection to renormalisation group in field theory;
overview of software: ITensor (Julia), TensorFlow Quantum, Quimb.

\medskip\textbf{Core learning outcomes:}
\begin{itemize}[leftmargin=1.4em,topsep=2pt,itemsep=1pt]
  \item Read and draw tensor network diagrams and identify contraction patterns.
  \item Distinguish MPS, TTN, MERA, and PEPS in terms of entanglement structure and cost.
  \item Explain tensor hypercontraction and its computational advantage.
  \item Connect tensor networks to renormalisation group and machine learning.
\end{itemize}

\medskip\textbf{Dual-track content:}
\begin{itemize}[leftmargin=1.4em,topsep=2pt,itemsep=1pt]
  \item \textit{Quantum chemistry/algorithms track:} THC for two-electron integrals;
    PEPS for 2D correlated materials; tensor-network quantum embedding; software workflows.
  \item \textit{Quantum information/mathematical track:} Entanglement structure of different TN
    geometries; holographic duality (AdS/MERA); MERA as a discrete realisation of AdS/CFT;
    complexity of TN contraction.
\end{itemize}

\medskip\textbf{Tier differentiation:}
HS: ``data stored as a network of multiplications rather than one giant array'' ---
tensor networks as efficient compression of high-dimensional data.
Beg.\ UG: matrix product of three matrices as a 3-tensor MPS; DMRG as 1D TN.
Adv.\ UG: MERA causal cone; PEPS contraction complexity; Tucker decomposition.
MSc: THC for electron repulsion integrals; Born machine; quantum-inspired ML.
PhD: AdS/MERA connection; PEPS for frustrated magnets; real-space RG from TN perspective.

\medskip\textbf{Assessment:}
Set~A --- draw and evaluate simple TN contractions (matrix chain, star network);
MPS representation of a GHZ state; MERA causal cone width.
Set~B --- derive SVD-based truncation error in TTN; THC factorisation of a $2\times2$ ERI block.
Projects: (simple) MPS compression of a random state: truncation error vs.\ bond dimension;
(complex) MERA ansatz for Ising critical point: scaling exponents from entanglement spectrum.
\end{tcolorbox}

\medskip

% ------ L09 -----------------------------------------------
\begin{tcolorbox}[modulebox, title={Module~II.1 --- Lecture L09: Natural Orbitals and the One-Particle Density Matrix}]
\textbf{90-min pacing:} (0--20~min) The one-particle reduced density matrix (1-RDM):
$\gamma(\mathbf{x};\mathbf{x}')=N\int\Psi^*(\mathbf{x}',\mathbf{x}_2,\ldots)\Psi(\mathbf{x},\mathbf{x}_2,\ldots)d\mathbf{x}_2\cdots d\mathbf{x}_N$;
one- and two-electron properties from the 1-RDM and 2-RDM;
$N$-representability of the 1-RDM: necessary conditions (eigenvalues in $[0,1]$, trace $=N$);
(20--45~min) Natural orbitals (NOs): eigenfunctions of the 1-RDM;
occupation numbers $n_i\in[0,1]$ (NOONs); $n_i=0$ or $1$ iff single Slater determinant;
fractional occupations as a measure of correlation; natural orbital expansion of $\Psi$
converges faster than canonical MO expansion; analogy with PCA / SVD of the wavefunction;
(45--65~min) Natural orbitals in practice: NOON-based active space selection;
iterative NO-based CI (iterative MRCI); CASSCF natural orbitals;
pair natural orbitals (PNO): biorthogonal pair functions defined from the 2-RDM;
PNO-CCSD(T) and DLPNO-CCSD(T): local correlation with PNO truncation;
(65--90~min) The 2-RDM and $N$-representability: P, Q, G conditions;
direct variational optimisation of the 2-RDM (2-RDM methods);
semidefinite programming (SDP) formulation; accuracy and scaling;
antisymmetric geminal power (AGP) and generalised valence bond (GVB) as
correlated pair wavefunctions.

\medskip\textbf{Core learning outcomes:}
\begin{itemize}[leftmargin=1.4em,topsep=2pt,itemsep=1pt]
  \item Define the 1-RDM and compute it for a two-electron wavefunction.
  \item Interpret NOONs as a measure of static correlation.
  \item Explain NOON-based active space selection and PNO local correlation.
  \item State the $N$-representability problem for the 2-RDM and the SDP approach.
\end{itemize}

\medskip\textbf{Dual-track content:}
\begin{itemize}[leftmargin=1.4em,topsep=2pt,itemsep=1pt]
  \item \textit{Computational chemistry track:} NOONs as diagnostics for multireference character;
    DLPNO-CCSD(T) active space from PNOs; GVB for bond breaking; SDP-2RDM software (Maple QMC).
  \item \textit{Mathematical/formal track:} 1-RDM as a Hermitian trace-class operator;
    Coleman's $N$-representability conditions; SDP duality; connection to quantum marginal problem.
\end{itemize}

\medskip\textbf{Tier differentiation:}
HS: ``occupation numbers tell you how many electrons are in each orbital'' ---
NOONs of 1 or 0 mean no correlation; numbers in between signal strong correlation.
Beg.\ UG: compute 1-RDM for He with two-configuration CI; diagonalise to get NOs; interpret NOONs.
Adv.\ UG: proof that NOs diagonalise the 1-RDM; NOON-based active space selection criterion.
MSc: PNO construction from 2-RDM; DLPNO truncation thresholds; $N$-representability P, Q, G conditions.
PhD: SDP for 2-RDM; Coleman's theorem; connection to quantum marginal problem; v2RDM.

\medskip\textbf{Assessment:}
Set~A --- compute NOONs for H$_2$ at stretched geometry; identify multireference character;
compare natural orbital expansion rate with canonical MO expansion.
Set~B --- prove that eigenfunctions of 1-RDM are NOs; derive the P condition for $N$-representability;
PNO truncation error bound.
Projects: (simple) NOON diagnostic: HF vs.\ FCI NOONs for H$_2$ dissociation;
(complex) 2-RDM SDP for H$_2$: comparison with FCI and CCSD.
\end{tcolorbox}

\medskip

% ------ L10 -----------------------------------------------
\begin{tcolorbox}[modulebox, title={Module~II.1 --- Lecture L10: Multi-Reference Methods}]
\textbf{90-min pacing:} (0--20~min) When single-reference fails: near-degeneracies,
bond breaking, transition metal chemistry, conical intersections;
static vs.\ dynamic correlation; multireference diagnostics: $T_1$ and $D_1$ amplitudes,
NOONs, $\langle\hat{S}^2\rangle$ deviation;
the complete active space (CAS) concept: partitioning MOs into inactive (doubly occupied),
active (variably occupied), and virtual (empty);
(20--45~min) CASSCF: simultaneous optimisation of CI coefficients and orbitals;
Newton--Raphson and DIIS for orbital rotations; super-CI and DIIS;
second-order CASSCF (CASSCF-2) and RASSCF (restricted active space);
active space selection heuristics: chemical intuition, NOON-based, automated (AVAS, DMET);
(45--70~min) MRCI: configuration expansion from CASSCF reference; MRCI+Q (Davidson correction)
for size-extensivity; NEVPT2 (strongly-contracted, partially-contracted):
second-order perturbation theory on a CASSCF reference; CASPT2 and RASPT2;
MRMP2 and QDPT; accuracy and applicability comparison;
(70--90~min) Automated active space methods: AutoCAS; AVAS (atomic valence active space);
embedding: DMET (density matrix embedding theory); SEET (self-energy embedding theory);
multireference methods in periodic systems: CASSCF for solids (preview);
MR-EOM-CC for excited states; summary and method choice guidelines.

\medskip\textbf{Core learning outcomes:}
\begin{itemize}[leftmargin=1.4em,topsep=2pt,itemsep=1pt]
  \item Identify situations where multireference methods are necessary using diagnostic criteria.
  \item Describe the CASSCF algorithm and the role of orbital optimisation.
  \item Compare MRCI+Q, NEVPT2, and CASPT2 in accuracy and computational cost.
  \item Apply automated active space selection (AVAS, DMET) to a chemical problem.
\end{itemize}

\medskip\textbf{Dual-track content:}
\begin{itemize}[leftmargin=1.4em,topsep=2pt,itemsep=1pt]
  \item \textit{Applied computational chemistry track:} Method selection for transition metals,
    photochemistry, and bond breaking; benchmark datasets (GMTKN55, MR-TDDFT comparison);
    active space size limitations and DMRG as the remedy.
  \item \textit{Formal/mathematical track:} MRCI as a truncated CI from a multi-determinant
    reference; Davidson correction and its size-extensivity restoration; NEVPT2 from
    second-order Rayleigh--Schr\"odinger PT on a CAS Hamiltonian.
\end{itemize}

\medskip\textbf{Tier differentiation:}
HS: ``when one configuration is not enough, use a committee'' ---
CASSCF as choosing the most important configurations from a controlled active space.
Beg.\ UG: minimal active space for H$_2$ bond breaking; CASSCF(2,2) energy curve.
Adv.\ UG: CASSCF orbital gradient and Hessian; Davidson correction derivation; NEVPT2 Hamiltonian.
MSc: RASSCF partitioning; AutoCAS/AVAS; DMET embedding; multireference diagnostics.
PhD: MRCI internal contraction; NEVPT2 strongly vs.\ partially contracted; MR-EOM-CC;
CASSCF for solids.

\medskip\textbf{Assessment:}
Set~A --- choose active space for O$_3$, Fe(CO)$_5$, and a Cr complex using NOONs;
CASSCF(2,2) energy for H$_2$; $T_1$ diagnostic interpretation.
Set~B --- derive Davidson correction; NEVPT2 second-order energy formula derivation;
CASSCF orbital gradient equation.
Projects: (simple) CASSCF vs.\ RHF potential energy surface for N$_2$ dissociation;
(complex) NEVPT2 vs.\ CASPT2 excitation energies for a transition metal complex.
\end{tcolorbox}

\medskip

% ------ L11 -----------------------------------------------
\begin{tcolorbox}[modulebox, title={Module~II.1 --- Lecture L11: Benchmark Systems and Accuracy Hierarchies}]
\textbf{90-min pacing:} (0--20~min) Philosophy of benchmarking: what does ``accurate'' mean?
Comparison to experiment vs.\ to higher-level theory; error cancellation; basis set vs.\
method error; systematic and non-systematic errors; error statistics: MAE, RMSE, MAX;
(20--45~min) Benchmark systems and datasets:
H$_2$ potential energy curve --- analytic benchmarks;
helium atom: exact (Pekeris) vs.\ HF vs.\ MP2 vs.\ CCSD(T) vs.\ FCI;
LiH: basis set convergence; N$_2$ dissociation: multireference character onset;
GMTKN55 composite benchmark: thermochemistry (W4-11), reaction barriers (BH76),
non-covalent interactions (S66, WATER27);
(45--65~min) Accuracy hierarchy:
HF $\to$ MP2 $\to$ CCSD $\to$ CCSD(T) $\to$ CCSDT $\to$ FCI;
scaling: $N^4$, $N^5$, $N^6$, $N^7$, $N^8$, exponential;
DFT comparison: relative accuracy and cost;
QMC in the hierarchy: DMC vs.\ CCSD(T) and FCI;
basis set convergence: CBS extrapolation schemes (Helgaker, Peterson);
composite methods: G3, W4, HEAT;
(65--90~min) Error analysis in practice: basis set incompleteness error (BSIE);
core--valence correlation; scalar relativistic effects; spin--orbit coupling;
diagonal Born--Oppenheimer correction; zero-point vibrational energy;
thermochemical accuracy: sub-kJ/mol goal; error budget for atomisation energies.

\medskip\textbf{Core learning outcomes:}
\begin{itemize}[leftmargin=1.4em,topsep=2pt,itemsep=1pt]
  \item Rank electronic structure methods by accuracy and computational scaling.
  \item Apply CBS extrapolation to estimate the basis set limit.
  \item Decompose the total error budget for a thermochemical property.
  \item Interpret benchmark statistics (MAE, RMSE, MAX) for method comparison.
\end{itemize}

\medskip\textbf{Dual-track content:}
\begin{itemize}[leftmargin=1.4em,topsep=2pt,itemsep=1pt]
  \item \textit{Applied track:} GMTKN55 and HEAT benchmarks as standards of chemical accuracy;
    composite methods for thermochemistry; relativistic and DBOC corrections; choosing
    method for a given system and accuracy target.
  \item \textit{Formal/statistical track:} Error propagation; CBS extrapolation theory;
    Helgaker $n^{-3}$ correlation energy extrapolation; systematic error cancellation
    in isodesmic reactions.
\end{itemize}

\medskip\textbf{Tier differentiation:}
HS: ``higher-level methods are more accurate but cost more time on the computer'' ---
accuracy ladder as a trade-off between exactness and feasibility.
Beg.\ UG: compare HF, MP2, CCSD(T) for H$_2$O geometry; CBS extrapolation for a simple molecule.
Adv.\ UG: full error budget for N$_2$ atomisation energy; W4 methodology components.
MSc: GMTKN55 analysis; relativistic corrections; DBOC evaluation; composite method design.
PhD: complete basis set limit theory; stochastic error in QMC benchmarks; renormalised MBPT.

\medskip\textbf{Assessment:}
Set~A --- rank methods by accuracy and cost; CBS extrapolation for He and H$_2$O;
compute MAE for a set of 5 atomisation energies.
Set~B --- derive Helgaker CBS formula; error budget for N$_2$ dissociation energy;
design a composite method for sub-kJ/mol thermochemistry.
Projects: (simple) accuracy ladder plot for H$_2$O: HF through CCSD(T)/CBS with error bars;
(complex) full W4-style error budget for C$_2$H$_2$: CCSD(T), core--valence, relativity, DBOC, ZPE.
\end{tcolorbox}

\medskip

% ------ L12 -----------------------------------------------
\begin{tcolorbox}[modulebox, title={Module~II.1 --- Lecture L12: Method Selection in Practice}]
\textbf{90-min pacing:} (0--20~min) The method selection problem: given a molecule, property,
and accuracy target, which method is optimal?
Decision tree: system size $N$; required accuracy; multireference character;
available computational resources; time constraints;
(20--45~min) Cost-scaling summary: HF $\mathcal{O}(N^3)$--$\mathcal{O}(N^4)$ (conventional);
MP2 $\mathcal{O}(N^5)$; CCSD $\mathcal{O}(N^6)$; CCSD(T) $\mathcal{O}(N^7)$;
CASSCF $\mathcal{O}(e^{N_{\rm act}})$; DMRG $\mathcal{O}(m^3 N)$; DMC $\mathcal{O}(N^3)$;
DFT $\mathcal{O}(N^3)$ (hybrid); linear-scaling DFT $\mathcal{O}(N)$;
crossover system sizes for practical use cases;
(45--65~min) Applicability regimes:
single-reference systems: CCSD(T)/CBS as gold standard; DFT for large single-reference;
multireference: CASSCF/NEVPT2 for moderate active spaces; DMRG for large active spaces;
solids and extended systems: DFT, periodic HF, embedding (DMET, QMMM);
excited states: EOM-CC, TDDFT, MR-EOM-CC, ADC;
non-equilibrium and dynamics: TDDFT, MCTDH, ML potentials;
(65--90~min) Practical workflow design: choosing basis set, method, software, and hardware;
black-box vs.\ expert use; automation and high-throughput screening;
machine learning potentials (MLPs): how DFT or CCSD(T) data trains MLPs;
future outlook: quantum computing approaches (VQE, QITE); fault-tolerant quantum chemistry.

\medskip\textbf{Core learning outcomes:}
\begin{itemize}[leftmargin=1.4em,topsep=2pt,itemsep=1pt]
  \item Apply the decision tree to select an appropriate method for a given system and target.
  \item Compare formal scaling with practical crossover sizes for common methods.
  \item Design a complete computational workflow from structure to property.
  \item Describe the role of ML potentials and quantum computing in the future of the field.
\end{itemize}

\medskip\textbf{Dual-track content:}
\begin{itemize}[leftmargin=1.4em,topsep=2pt,itemsep=1pt]
  \item \textit{Practical/industrial track:} High-throughput screening; automation with AiiDA
    or ASE; ORCA and MOLPRO workflows; ML potential training pipelines; QM/MM for enzymes.
  \item \textit{Foundational track:} Formal complexity and why crossovers occur;
    locality of correlation; quantum computing resources for quantum chemistry;
    resource estimation for fault-tolerant quantum chemistry.
\end{itemize}

\medskip\textbf{Tier differentiation:}
HS: ``match the tool to the job'' --- a pocket guide to choosing a method;
analogy with choosing the right level of map detail for your journey.
Beg.\ UG: method decision tree for 5 representative molecules; identify bottleneck.
Adv.\ UG: crossover analysis; formal scaling vs.\ prefactor discussion; workflow design.
MSc: high-throughput screening design; ML potential training workflow; embedding strategy.
PhD: quantum computing resource estimation for FeMoco; fault-tolerant vs.\ NISQ;
future of the field: where will the next major accuracy gains come from?

\medskip\textbf{Assessment:}
Set~A --- method selection for: (i) drug-like molecule geometry, (ii) transition metal
  reaction barrier, (iii) solid-state band gap, (iv) photochemical conical intersection.
Set~B --- derive crossover size between CCSD(T) and DMC; design an error budget and workflow
for sub-kJ/mol prediction of the N$_2$ dissociation energy.
Projects: (simple) decision-tree web app for method selection;
(complex) full computational study of a reaction: method convergence, basis set extrapolation,
and comparison to experiment.
\end{tcolorbox}

%------------------------------------------------------------
\subsubsection*{Assessment Strategy}
\addcontentsline{toc}{subsubsection}{Module II.1 Assessment Strategy}

\begin{tcolorbox}[highlightbox]
\textbf{Per-Lecture Assessment Bundle} (every lecture):
concept check (10--15~min, 5--8 short conceptual questions);
Problem Set~A (simple, 6--10 computational exercises);
Problem Set~B (complex, 3--6 multi-step derivations and proofs);
mini-project simple (1--3 hours);
mini-project complex (5--10 hours, open-ended);
research questions well-defined (2--4 scoped prompts);
research questions fully open (1--2 frontier-style prompts).
\end{tcolorbox}

\medskip
\noindent\textbf{Course-level grading template:}
\begin{itemize}[leftmargin=1.5em]
  \item Weekly homework (best 10/12): 40\%
  \item Method comparison project (after L06): 20\%
  \item Computational research project: 25\%
  \item Final exam (take-home or whiteboard): 15\%
\end{itemize}
Mixed cohorts: Beg.\ UG $\to$ Set~A + simple projects; MSc/PhD $\to$ Set~B + complex projects.

%------------------------------------------------------------
\subsubsection*{Reading Strategy}
\addcontentsline{toc}{subsubsection}{Module II.1 Reading Strategy}

\noindent\textbf{Track A --- Core textbooks:}
Helgaker, J\o rgensen \& Olsen, \textit{Molecular Electronic-Structure Theory} (Wiley, 2000) ---
the definitive reference for Chapters 3--14;
Szabo \& Ostlund, \textit{Modern Quantum Chemistry} (Dover) --- HF and post-HF primer.

\noindent\textbf{Track B --- Methods references:}
Jensen, \textit{Introduction to Computational Chemistry} (Wiley) --- accessible survey;
Cramer, \textit{Essentials of Computational Chemistry} (Wiley);
Nightingale \& Umrigar (eds.), \textit{Quantum Monte Carlo Methods in Physics and Chemistry} (Kluwer).

\noindent\textbf{Track C --- Advanced and frontier:}
Schollw\"ock, Rev.\ Mod.\ Phys.\ \textbf{77} (2005) and \textbf{82} (2011) --- DMRG reviews;
Or\'us, Ann.\ Phys.\ \textbf{349} (2014) --- tensor networks;
Needs \textit{et al.}, J.\ Chem.\ Phys.\ \textbf{152} (2020) --- CASINO DMC review.

%=============================================================
\subsection{Module II.2 --- Path Integrals in Quantum Mechanics}
\label{sec:mod-II2}
%=============================================================

\begin{tcolorbox}[modulebox, title={Module~II.2 --- Lecture Overview}]
\begin{longtable}{clp{7cm}}
\toprule
\textbf{Lec.} & \textbf{Title} & \textbf{Content}\\
\midrule
\endhead
L01 & Feynman's Postulate & Sum over histories; Lagrangian action; classical limit\\
L02 & Derivation of the Path Integral & Trotter factorisation; discretisation; continuum limit\\
L03 & Gaussian Path Integrals & Harmonic oscillator; stationary phase; fluctuation determinant\\
L04 & Saddle Point and WKB & Semiclassical expansion; tunnelling amplitude\\
L05 & Imaginary-Time Path Integrals & Wick rotation; partition function; Euclidean action\\
L06 & Statistical Mechanics Connection & $Z=\operatorname{Tr}e^{-\beta\hat{H}}$ as path integral\\
L07 & Coherent State Path Integral & Bosonic/fermionic coherent states; Grassmann variables\\
L08 & Instantons and Tunnelling & Bounce solution; decay rate; double-well potential\\
L09 & Path Integrals in Field Theory & $\phi^4$ theory; Feynman rules; propagator\\
L10 & Gauge Fields & Aharonov--Bohm; minimal coupling; Wilson loop\\
L11 & Numerical Path Integrals & Monte Carlo; path integral MD\\
L12 & Berry Phase from Path Integrals & Geometric action; holonomy; applications\\
\bottomrule
\end{longtable}
\end{tcolorbox}

%=============================================================
\subsection{Module II.3 --- Theory of Angular Momentum}
\label{sec:mod-II3}
%=============================================================

\begin{tcolorbox}[modulebox, title={Module~II.3 --- Lecture Overview}]
\begin{longtable}{clp{7cm}}
\toprule
\textbf{Lec.} & \textbf{Title} & \textbf{Content}\\
\midrule
\endhead
L01 & $\mathfrak{su}(2)$ Algebra & $[\hat{J}_i,\hat{J}_j]=i\hbar\epsilon_{ijk}\hat{J}_k$; Casimir $\hat{\mathbf{J}}^2$\\
L02 & Ladder Operators and Spectrum & $\hat{J}_\pm$; $\ket{j,m}$; quantisation of $j$\\
L03 & Orbital Angular Momentum & $\hat{\mathbf{L}}=-i\hbar(\mathbf{r}\times\nabla)$; spherical harmonics\\
L04 & Spin Angular Momentum & Pauli matrices; spinors; $SU(2)$ double cover of $SO(3)$\\
L05 & Addition of Angular Momenta & Clebsch--Gordan coefficients; $j_1\otimes j_2$\\
L06 & Wigner $3j$-Symbols and Racah Algebra & Selection rules; Wigner--Eckart theorem\\
L07 & Spherical Tensor Operators & Rank-$k$ tensors; matrix element theorems\\
L08 & Hydrogen Atom Full Solution & $SO(4)$ symmetry; $n,l,m,m_s$ quantum numbers\\
L09 & Fine and Hyperfine Structure & Relativistic corrections; spin--orbit; contact term\\
L10 & Wigner $D$-matrices & $D^{(j)}_{mm'}(\alpha,\beta,\gamma)$; applications\\
L11 & Angular Momentum in Many-Body Systems & Shell model; $LS$ and $jj$ coupling\\
L12 & Rotational Spectroscopy Connection & Rigid rotor; selection rules; molecular spectra\\
\bottomrule
\end{longtable}
\end{tcolorbox}

%=============================================================
\subsection{Module II.4 --- Second Quantisation and Introductory Many-Body Theory}
\label{sec:mod-II4}
%=============================================================

\begin{tcolorbox}[modulebox, title={Module~II.4 --- Lecture Overview}]
\begin{longtable}{clp{7cm}}
\toprule
\textbf{Lec.} & \textbf{Title} & \textbf{Content}\\
\midrule
\endhead
L01 & Fock Space & Occupation number representation; bosons and fermions;
     $[\hat{a}_i,\hat{a}_j^\dagger]=\delta_{ij}$ vs.\ $\{\hat{c}_i,\hat{c}_j^\dagger\}=\delta_{ij}$\\
L02 & Operators in Second Quantisation & One- and two-body operators; normal ordering; Wick's theorem\\
L03 & Homogeneous Electron Gas & Jellium model; Fermi sea; exchange energy\\
L04 & Hartree--Fock in Second Quantisation & Mean-field decoupling; HF equations; Koopman's theorem\\
L05 & Beyond Mean Field & Configuration interaction; MP2; diagrammatics\\
L06 & Green's Functions I & Single-particle GF; Lehmann representation; spectral function\\
L07 & Green's Functions II & Dyson equation; self-energy; quasiparticles\\
L08 & Interacting Bosons: BEC & Bogoliubov transformation; phonon spectrum; superfluidity\\
L09 & Interacting Fermions: BCS & Cooper pairs; gap equation; Bogoliubov--Valatin transformation\\
L10 & The Hubbard Model & Tight-binding + interactions; half-filling; Mott transition\\
L11 & Quantum Magnetism & Heisenberg exchange; spin waves; Holstein--Primakoff\\
L12 & Towards Non-Equilibrium & Keldysh formalism motivation; drive-dissipation balance\\
\bottomrule
\end{longtable}
\end{tcolorbox}

%=============================================================
\subsection{Module II.5 --- Non-Equilibrium Quantum Many-Body Green's Functions}
\label{sec:mod-II5}
%=============================================================

\begin{tcolorbox}[modulebox, title={Module~II.5 --- Lecture Overview}]
\begin{longtable}{clp{7cm}}
\toprule
\textbf{Lec.} & \textbf{Title} & \textbf{Content}\\
\midrule
\endhead
L01 & Motivation and Scope & Transport; ultrafast dynamics; driven systems; failures of equilibrium\\
L02 & The Keldysh Contour & Closed time path; forward/backward branches; contour ordering\\
L03 & Contour-Ordered Green's Functions & $G^c$, $G^<$, $G^>$, $G^R$, $G^A$, $G^K$\\
L04 & Langreth Rules & Analytic continuation; convolution on contour\\
L05 & Kadanoff--Baym Equations & Functional derivation; two-time structure; initial correlations\\
L06 & Self-Energies and Approximations & Hartree--Fock; GW; second Born; T-matrix\\
L07 & Quantum Boltzmann Equation & Gradient expansion; kinetic limit; quasi-classical transport\\
L08 & Steady-State Transport & Meir--Wingreen formula; Landauer--B\"uttiker; molecular junctions\\
L09 & Electron--Phonon Coupling & Fr\"ohlich Hamiltonian; polaron; phonon self-energy\\
L10 & Pump--Probe Dynamics & Two-time spectra; transient absorption; non-equilibrium spectral functions\\
L11 & Driven-Dissipative Systems & Lindblad from Keldysh; NEGF + bath; Floquet--Keldysh\\
L12 & Numerical Methods & Time-stepping KB equations; HPC considerations; benchmarks\\
\bottomrule
\end{longtable}
\end{tcolorbox}

%=============================================================
\subsection{Module II.6 --- Light--Matter Interaction, Linear Response, and Spectroscopy}
\label{sec:mod-II6}
%=============================================================

\begin{tcolorbox}[modulebox, title={Module~II.6 --- Lecture Overview}]
\begin{longtable}{clp{7cm}}
\toprule
\textbf{Lec.} & \textbf{Title} & \textbf{Content}\\
\midrule
\endhead
L01 & Classical and Quantum EM Fields & Vector potential; Coulomb gauge; quantisation; photon\\
L02 & Light--Matter Coupling & Minimal coupling; dipole approximation; PZW transformation\\
L03 & Linear Response Theory & Kubo formula; retarded response $\chi^R(\omega)$; causality\\
L04 & Fluctuation--Dissipation Theorem & FDT; fluctuations from dissipation; thermal equilibrium\\
L05 & Optical Conductivity & Drude model from Kubo; $f$-sum rule; Kramers--Kronig relations\\
L06 & Absorption and Emission & Fermi golden rule for radiation; Einstein $A$ and $B$ coefficients\\
L07 & Selection Rules & Dipole, quadrupole; parity; angular momentum conservation; Laporte rule\\
L08 & Non-Linear Response & $\chi^{(2)},\chi^{(3)}$; two-photon absorption; harmonic generation\\
L09 & Time-Domain Spectroscopy & Free-induction decay; response from correlation function\\
L10 & Raman and Resonance Raman & Vibronic coupling; Kramers--Heisenberg formula\\
L11 & Photoelectron Spectroscopy & XPS/UPS; sudden approximation; Dyson orbital; shake-up\\
L12 & Cavity QED and Strong Coupling & Jaynes--Cummings; vacuum Rabi splitting; polaritons\\
\bottomrule
\end{longtable}
\end{tcolorbox}

%=============================================================
\subsection{Module II.7 --- Quantum Statistical Mechanics in Equilibrium}
\label{sec:mod-II7}
%=============================================================

\begin{tcolorbox}[modulebox, title={Module~II.7 --- Lecture Overview}]
\begin{longtable}{clp{7cm}}
\toprule
\textbf{Lec.} & \textbf{Title} & \textbf{Content}\\
\midrule
\endhead
L01 & Density Operators in Thermodynamics & Gibbs state $\hat{\rho}=e^{-\beta\hat{H}}/Z$; entropy\\
L02 & Canonical and Grand-Canonical Ensembles & Partition function; chemical potential; grand potential\\
L03 & Quantum Ideal Gases & Bose--Einstein and Fermi--Dirac statistics; degeneracy pressure\\
L04 & Bose--Einstein Condensation & Critical temperature; condensate fraction; superfluid helium\\
L05 & Fermi Liquid Theory & Quasiparticles; Landau parameters; specific heat; compressibility\\
L06 & Imaginary-Time Formalism & Matsubara Green's functions; fermionic/bosonic frequencies $\omega_n$\\
L07 & Perturbation Theory at Finite $T$ & Diagrammatics; Dyson equation; self-energy at $T>0$\\
L08 & Phase Transitions I & Landau theory; order parameter; spontaneous symmetry breaking\\
L09 & Phase Transitions II & Ising model; mean-field; universality; quantum phase transitions\\
L10 & Quantum Criticality & Quantum critical point; $\omega/T$ scaling; non-Fermi liquid\\
L11 & Entanglement in Many-Body Systems & Area vs.\ volume law; topological entropy\\
L12 & Open Problems and Frontiers & High-$T_c$ superconductivity; fractional QH; quantum simulation\\
\bottomrule
\end{longtable}
\end{tcolorbox}

\newpage

%#############################################################
%  SERIES III SEPARATOR PAGE
%#############################################################
\newpage
\thispagestyle{empty}
\pagecolor{darkblue}
\color{white}
\vspace*{2.2cm}
\begin{center}
  {\color{accentgold}\rule{11cm}{1.5pt}}\\[0.5cm]
  {\fontsize{11}{14}\selectfont\bfseries\color{accentgold}\MakeUppercase{Series III}}\\[0.3cm]
  {\fontsize{22}{28}\selectfont\bfseries Molecular Electronic Structure}\\[0.2cm]
  {\fontsize{22}{28}\selectfont\bfseries \& Quantum Chemistry}\\[0.4cm]
  {\color{accentgold}\rule{11cm}{1.5pt}}\\[1.0cm]
  {\large\itshape Three modules --- 31 lectures $\times$ 90 minutes}\\[0.5cm]
  {\normalsize\itshape Audience: Graduate students in chemistry and chemical physics}\\[0.3cm]
  {\normalsize\itshape Prerequisites: Series~I required; Modules II.1 and II.4 strongly recommended}\\[1.6cm]

  \begin{tabular}{@{}clr@{}}
    \toprule
    \textbf{Module} & \textbf{Title} & \textbf{Lectures}\\
    \midrule
    \textbf{III.1} & Basic Premises of Quantum Chemistry & 3\\[3pt]
                   & \multicolumn{2}{@{}l}{\small\color{accentgold}%
                     Molecular Hamiltonian ($\hat{T}_n+\hat{T}_e+\hat{V}_{ne}+\hat{V}_{ee}+\hat{V}_{nn}$);}\\
                   & \multicolumn{2}{@{}l}{\small\color{accentgold}%
                     Born--Oppenheimer approximation; adiabatic potential energy surfaces;}\\
                   & \multicolumn{2}{@{}l}{\small\color{accentgold}%
                     non-adiabatic couplings; Feynman--Hellmann theorem; force calculations.}\\[6pt]
    \textbf{III.2} & Principles of Molecular and Atomic Electronic Structure & 16\\[3pt]
                   & \multicolumn{2}{@{}l}{\small\color{accentgold}%
                     Hydrogenic systems; polyelectronic wavefunctions; spin orbitals;}\\
                   & \multicolumn{2}{@{}l}{\small\color{accentgold}%
                     Slater determinants; Hartree--Fock; Roothaan--Hall; Koopman's theorem;}\\
                   & \multicolumn{2}{@{}l}{\small\color{accentgold}%
                     configuration interaction; orbital theories; natural and localised orbitals;}\\
                   & \multicolumn{2}{@{}l}{\small\color{accentgold}%
                     basis sets (STO, GTO, Pople, Dunning, Karlsruhe); molecular integrals.}\\[6pt]
    \textbf{III.3} & Density Functional Theory and Beyond & 12\\[3pt]
                   & \multicolumn{2}{@{}l}{\small\color{accentgold}%
                     Hohenberg--Kohn theorems; Thomas--Fermi--Weizs\"acker; Kohn--Sham DFT;}\\
                   & \multicolumn{2}{@{}l}{\small\color{accentgold}%
                     XC functionals (LDA, GGA, hybrid, double-hybrid); Fukui functions;}\\
                   & \multicolumn{2}{@{}l}{\small\color{accentgold}%
                     tight-binding; H\"uckel/M\"obius; crystal/ligand field theory;}\\
                   & \multicolumn{2}{@{}l}{\small\color{accentgold}%
                     full CI QMC; diffusion QMC; method comparison and outlook.}\\[4pt]
    \midrule
    & \textbf{Total} & \textbf{31}\\
    \bottomrule
  \end{tabular}

  \vfill
  {\color{accentgold}\rule{11cm}{0.8pt}}\\[0.4cm]
  {\small\itshape\color{white!80!darkblue}%
    Focus: first-principles molecular physics $\cdot$
    From exact wavefunction methods to DFT $\cdot$
    Computational \& theoretical depth}\\[0.3cm]
  {\color{accentgold}\rule{11cm}{0.8pt}}
\end{center}
\newpage
\nopagecolor
\color{black}

%#############################################################
\section{Series III --- Molecular Electronic Structure \& Quantum Chemistry}
\label{sec:seriesIII}
%#############################################################

\begin{tcolorbox}[seriesbox, title={Series~III Overview}]
\textbf{Prerequisite:} Series~I required; Series~II Modules II.1, II.4 strongly recommended.\\
\textbf{Structure:} Three modules totalling approximately 31 lectures.\\
\textbf{Core aim:} Rigorous, graduate-level understanding of molecular electronic structure
from first principles through to modern computational methods.\\[4pt]
\begin{tabular}{llr}
\toprule
\textbf{Module} & \textbf{Title} & \textbf{Lectures}\\
\midrule
III.1 & Basic Premises of Quantum Chemistry & 3\\
III.2 & Principles of Molecular and Atomic Electronic Structure & 16\\
III.3 & Density Functional Theory and Beyond & 12\\
\midrule
 & \textbf{Total} & \textbf{31}\\
\bottomrule
\end{tabular}
\end{tcolorbox}

%=============================================================
\subsection{Module III.1 --- Basic Premises of Quantum Chemistry}
\label{sec:mod-III1}
%=============================================================

\begin{tcolorbox}[chembox, title={Module~III.1 --- Lecture Overview}]
\begin{longtable}{clp{8cm}}
\toprule
\textbf{Lec.} & \textbf{Title} & \textbf{Content}\\
\midrule
\endhead
L01 & The Molecular Hamiltonian &
     $\hat{H}=\hat{T}_n+\hat{T}_e+\hat{V}_{ne}+\hat{V}_{ee}+\hat{V}_{nn}$;
     atomic units; degrees of freedom\\
L02 & Born--Oppenheimer Approximation and Adiabaticity &
     Electronic vs.\ nuclear time scales; adiabatic PES;
     non-adiabatic couplings; derivative couplings\\
L03 & Parametric Dependence and the Feynman--Hellmann Theorem &
     $\partial_\lambda E = \langle\partial_\lambda\hat{H}\rangle$;
     force calculations; geometry optimisation; potential energy gradients\\
\bottomrule
\end{longtable}
\end{tcolorbox}

%=============================================================
\subsection{Module III.2 --- Principles of Molecular and Atomic Electronic Structure}
\label{sec:mod-III2}
%=============================================================

\begin{tcolorbox}[chembox, title={Module~III.2 --- Lecture Overview (Electronic Hamiltonian)}]
\begin{longtable}{clp{7.2cm}}
\toprule
\textbf{Lec.} & \textbf{Title} & \textbf{Content}\\
\midrule
\endhead
L01 & The Quantum Many-Body Problem in Chemistry &
      Complexity of $\hat{H}_e$; $N$-representability; curse of dimensionality\\
L02 & Hydrogenic Systems &
      Exact solution; $Z_{\rm eff}$; quantum defect; one-electron atomic orbitals; radial nodes\\
L03 & Polyelectronic Systems &
      Helium as prototype; electron correlation; role of inter-electronic repulsion\\
L04 & Polyelectronic Wavefunctions, Spin Orbitals, and Slater Screening &
      Spin orbital $\chi(\mathbf{x})=\phi(\mathbf{r})\sigma(s)$;
      Slater rules for $Z_{\rm eff}$; shielding constants\\
L05 & Antisymmetry, Hartree Products, and Slater Determinants &
      Pauli exclusion; antisymmetric wavefunction;
      Slater determinant $|\chi_1\cdots\chi_N|$\\
L06 & Restricted and Unrestricted Methods &
      RHF vs.\ UHF; spatial orbital constraints;
      spin contamination $\langle\hat{S}^2\rangle$\\
L07 & Variational Methods &
      Basic Ritz; Hylleraas--Mayer canonical form;
      Hylleraas--Undheim theorem; bracketing theorem\\
L08 & Electronic Configurations, Aufbau, and Virtual Orbitals &
      Orbital energy ordering; general Aufbau principle;
      virtual orbitals; excited configurations\\
L09 & Configuration Interaction and the Exact Non-Relativistic Limit &
      Full CI expansion; truncated CI (CISD, CISDT);
      exact non-relativistic ground state; size consistency\\
L10 & Hartree--Fock Theory, Koopman's Theorem, and Correlation Energy &
      SCF procedure; Roothaan--Hall equations;
      Koopman's theorem for IPs/EAs;
      $E_{\rm corr}=E_{\rm exact}-E_{\rm HF}$\\
L11 & Orbital Theories and Bonding Models &
      MO theory; Valence Bond theory;
      Coulson--Fischer wavefunction; resonance structures\\
L12 & Natural and Localised Orbitals &
      Natural Bond Orbitals (NBO); L\"owdin orbitals;
      Natural Atomic Orbitals (NAO);
      Weinhold approximation; Foster--Boys localisation\\
L13 & Basis Sets I --- Types of Functions &
      Slater-type orbitals (STO); Gaussian-type orbitals (GTO);
      contraction; minimal, split-valence, polarised\\
L14 & Basis Sets II --- Standard Basis Families &
      Pople (6-31G*); Dunning cc-pVXZ; Karlsruhe def2;
      frozen-core approximation; effective core potentials (ECP)\\
L15 & Molecular Integral Schemes &
      Two-electron integrals; Obara--Saika recurrence;
      Huzinaga auxiliary; Helgaker density fitting; RI approximation\\
L16 & Glimpse into Solid State &
      PAW approximation; plane-wave basis;
      contrast with GTO; Bloch electrons revisited;
      band structure vs.\ molecular orbital picture\\
\bottomrule
\end{longtable}
\end{tcolorbox}

%=============================================================
\subsection{Module III.3 --- Density Functional Theory and Beyond}
\label{sec:mod-III3}
%=============================================================

\begin{tcolorbox}[chembox, title={Module~III.3 --- Lecture Overview (DFT, Effective Models, and Special Methods)}]
\begin{longtable}{clp{7.2cm}}
\toprule
\textbf{Lec.} & \textbf{Title} & \textbf{Content}\\
\midrule
\endhead
L01 & Hohenberg--Kohn Theorems &
      First theorem ($\rho\to V_{\rm ext}$);
      second theorem (variational); $v$-representability\\
L02 & Thomas--Fermi--Weizs\"acker Theory &
      Orbital-free DFT; TF kinetic energy;
      gradient correction; limitations\\
L03 & Kohn--Sham DFT &
      KS equations; non-interacting reference;
      XC functional; self-consistent field\\
L04 & Overview of XC Functionals &
      Jacob's ladder; LDA, GGA (PBE), meta-GGA,
      hybrid (B3LYP, PBE0), range-separated, double-hybrid; Perdew's ladder\\
L05 & Useful DFT Quantities &
      Fukui functions $f^+,f^-,f^0$;
      Pearson hardness/softness;
      Yang--Parr electronegativity $\chi=-\mu$; dual descriptor\\
L06 & Tight-Binding and Effective Hamiltonians &
      LCAO tight-binding; Hamiltonian matrix elements;
      Slater--Koster parameters; SSH model\\
L07 & H\"uckel and M\"obius Aromaticity Models &
      Simple H\"uckel MO theory;
      $4n$ vs.\ $4n+2$ H\"uckel rule;
      M\"obius topology; Woodward--Hoffmann rules\\
L08 & Crystal Field Theory &
      $d$-orbital splitting in octahedral, tetrahedral;
      $\Delta_o$; high-spin/low-spin; Jahn--Teller effect\\
L09 & Ligand Field Theory &
      Nephelauxetic effect; angular overlap model;
      Tanabe--Sugano diagrams;
      Racah--Van Vleck parameters $A,B,C$\\
L10 & Full CI Quantum Monte Carlo &
      Booth--Alavi FCIQMC algorithm;
      walker dynamics; sign problem in FCI space;
      initiator approximation; benchmark applications\\
L11 & Diffusion Quantum Monte Carlo &
      Branching random walk; fixed-node approximation;
      trial wavefunction quality; pseudopotentials in DMC\\
L12 & Synthesis and Outlook &
      Method comparison (accuracy, scaling, cost);
      current challenges: strong correlation,
      excited states, relativistic effects, periodic systems\\
\bottomrule
\end{longtable}
\end{tcolorbox}

\newpage
%=============================================================
\section*{Appendix A --- Programme Roadmap}
\addcontentsline{toc}{section}{Appendix A --- Programme Roadmap}
%=============================================================

\begin{center}
\begin{tabular}{@{}llll@{}}
\toprule
\textbf{Module} & \textbf{Title (abbreviated)} & \textbf{Core Prerequisites} & \textbf{Lectures}\\
\midrule
\textbf{I.1} & State Space, Operators, Bra-Ket & Linear algebra, classical mechanics & 6\\
\textbf{I.2} & Bra-Ket Notation (Dirac + Sakurai) & I.1 (concurrent or prior) & 12\\
\textbf{I.3} & Dynamics, Perturbation Theory, Systems & I.1; I.2 recommended & 11\\
\midrule
\textbf{II.1} & Ground-State Methods & Series~I & 12\\
\textbf{II.2} & Path Integrals & Series~I & 12\\
\textbf{II.3} & Angular Momentum & Series~I & 12\\
\textbf{II.4} & Second Quantisation & Series~I, II.3 & 12\\
\textbf{II.5} & NEGF & II.4, II.7 & 12\\
\textbf{II.6} & Light-Matter/Spectroscopy & Series~I, II.4 & 12\\
\textbf{II.7} & Quantum Stat.\ Mech. & Series~I & 12\\
\midrule
\textbf{III.1} & Basic Premises & Series~I & 3\\
\textbf{III.2} & Electronic Structure & Series~I, II.4 (rec.) & 16\\
\textbf{III.3} & DFT and Beyond & III.1, III.2 & 12\\
\midrule
 & \textbf{Grand total} & & \textbf{144}\\
\bottomrule
\end{tabular}
\end{center}

%=============================================================
\section*{Appendix B --- Suggested Textbooks}
\addcontentsline{toc}{section}{Appendix B --- Suggested Textbooks}
%=============================================================

\medskip\noindent{\large\bfseries\color{midblue}Quantum Mechanics --- Foundations}\par\smallskip
\begin{itemize}[leftmargin=1.5em]
  \item Dirac, P.A.M.\ --- \textit{The Principles of Quantum Mechanics}, 4th ed.\ (Oxford)
  \item Sakurai, J.J.\ \& Napolitano, J.\ --- \textit{Modern Quantum Mechanics}, 3rd ed.
  \item Weinberg, S.\ --- \textit{Lectures on Quantum Mechanics}, 2nd ed.
  \item Galindo, A.\ \& Pascual, P.\ --- \textit{Quantum Mechanics} (2 vols., Springer)
\end{itemize}

\medskip\noindent{\large\bfseries\color{midblue}Many-Body and Advanced Topics}\par\smallskip
\begin{itemize}[leftmargin=1.5em]
  \item Fetter, A.L.\ \& Walecka, J.D.\ --- \textit{Quantum Theory of Many-Particle Systems}
  \item Bruus, H.\ \& Flensberg, K.\ --- \textit{Introduction to Many-Body Quantum Theory in Condensed Matter Physics}
  \item Stefanucci, G.\ \& van Leeuwen, R.\ --- \textit{Nonequilibrium Many-Body Theory of Quantum Systems} (Cambridge)
  \item Negele, J.W.\ \& Orland, H.\ --- \textit{Quantum Many-Particle Systems}
\end{itemize}

\medskip\noindent{\large\bfseries\color{midblue}Quantum Chemistry and Electronic Structure}\par\smallskip
\begin{itemize}[leftmargin=1.5em]
  \item Szabo, A.\ \& Ostlund, N.S.\ --- \textit{Modern Quantum Chemistry}
  \item Helgaker, T., Jorgensen, P.\ \& Olsen, J.\ --- \textit{Molecular Electronic-Structure Theory} (Wiley)
  \item Jensen, F.\ --- \textit{Introduction to Computational Chemistry}, 3rd ed.
  \item Koch, W.\ \& Holthausen, M.C.\ --- \textit{A Chemist's Guide to Density Functional Theory}
  \item Martin, R.M., Reining, L.\ \& Ceperley, D.\ --- \textit{Interacting Electrons} (Cambridge)
\end{itemize}

%=============================================================
\end{document}
